\section{Agentic Image Generation, Editing, and Restoration}
\label{sec:image-generation}

Image creation agents must control an evolving artifact whose relevant state can include composition, masks, object identity, references, edit history, degradation estimates, user preferences, and tool outputs. Section~\ref{sec:analytical-framework} identifies the analytical components used to compare these systems. This section instead organizes image-domain studies by the primary contribution that they demonstrate. Each paper is assigned one location; secondary mechanisms remain relevant for cross-cutting comparison but do not determine a second placement.

\subsection{Feedback-Driven Image Generation and Editing}
Feedback-driven systems observe a generated artifact, identify a deviation from the current goal, and revise a later generation or editing decision. The following subsections distinguish compositional and design refinement from image-level editing, completion, and restoration.

\subsubsection{Compositional Generation and Design Refinement}
These systems use generated images as evidence for revising compositional constraints, design intent, or local defects. Their common contribution is a feedback action that changes the current generation trajectory.

\paragraph{Improving Compositional Text-to-Image Generation with Large Vision-Language Models.}
This work uses an LVLM as both evaluator and editor for compositional text-to-image generation. The LVLM converts prompts into QA checks, uses answer accuracy as a loss weight for ReFL fine-tuning, and at inference iteratively detects mismatches and directs SAM+Blended Diffusion to correct them until fully compliant. ~\cite{wen2023compositional}

\paragraph{Idea2Img.}
 is a GPT-4V agent that iteratively refines prompts for multimodal image design without prior knowledge of the T2I model's optimal usage. Each round, GPT-4V generates candidate prompts, selects the best draft, reflects on mismatches with the user IDEA (interleaved text and images), and stores feedback in memory to guide revision. On 104 queries, iterative SDXL v1.0 wins 56.7\% (vs. 29.8\% initial, 13.5\% manual; +26.9 SDXL, +16.3 IF). Memory guides revision within runs; cross-run retention is untested.~\cite{yang2024idea2img}

\paragraph{MuLan.}
 is a training-free agent that decomposes complex prompts into object-level subtasks, generating each object conditioned on prior ones via LLM-planned masks and attention guidance, with a VLM verifying each subtask and triggering regeneration on violation; users may intervene mid-process. MuLan outperforms vanilla SD v1.4 on 200 prompts (77.24\%/75.15\% vs. 66.43\%/56.73\%); VLM feedback proves critical.~\cite{li2025mulan}

\paragraph{VisualPrompter.}
 fixes semantic omissions in T2I generation by decomposing prompts into atomic concepts, using a VLM to identify missing concepts in the generated image, and revising only those units. Across DSG and TIFA benchmarks with four generators, it raises the average semantic score from 78.4 to 83.0.~\cite{wu2026visualprompter}

\paragraph{Maestro.}
 is a test-time agent system that automates T2I prompt refinement from underspecified user prompts. It decomposes prompts into DVQs, applies MLLM-driven targeted revisions with semantic verification, and maintains the best image via pairwise MLLM tournaments. Maestro outperforms the best baseline on both benchmarks (p2-hard: 0.921,   DSG-1K:0.882) .~\cite{avg250910704}

\paragraph{CRAFT}
 is a training-free framework that improves T2I generation via constraint verification. It converts prompts into visual questions, uses a VLM to identify failures with rationales, and applies targeted LLM edits only where needed, retaining the best image and stopping once all constraints pass.It raises VQA/DSG across backbones (e.g., FLUX-Schnell: 0.78→0.86/0.786→0.857)~\cite{kovalev2025craft}

\paragraph{Iterative Refinement Improves Compositional Image Generation.}
This work proposes iterative refinement for compositional T2I generation, where a VLM critic inspects the current image and selects among continue, restart, backtrack, or stop under a fixed budget, optionally combined with parallel sampling. On ConceptMix (k=7), It improves accuracy by 16.9\% over parallel sampling (13.8\% on T2I-CompBench 3D-spatial).

\paragraph{Agentic Flow Steering.}
AFS-Search is a training-free closed-loop framework on FLUX.1-dev for spatially grounded T2I generation. It rewrites prompts into explicit constraints, diagnoses intermediate latents, and compares three rollouts (baseline, exploration, and corrective via CLIP+SAM3 velocity modulation), selecting the best by VLM score. On T2I-CompBench, Pro scores 0.6748 (spatial: 0.6250); Fast scores 0.6219 at 32.5s. ~\cite{avg260318627}

\paragraph{FiRe.}
 is a fine-grained reasoning framework for T2I generation that decomposes prompts into atomic tuples, verifies each via VQA, and applies localized corrections only where needed; FiRe-GRPO supplies step-level rewards. ~\cite{kim2026fire}

\paragraph{Self-Reasoning Agentic Framework for Narrative Product Grid-Collage Generation.}
This framework generates product-advertising collages by decoupling narrative reasoning from pixel synthesis: agents plan a Product Narrative Framework (identity/usage/context/consumer), specify photographic decisions, and jointly synthesize the grid; a two-gate critique triggers targeted replanning or refinement on failure.~\cite{avg260416958}

\subsubsection{Iterative Image Editing, Completion, and Restoration}
This group places the feedback loop directly on an existing image. The controller diagnoses a mismatch, completion gap, or localized defect and selects a later editing or restoration operation.

\paragraph{An LLM--LVLM Driven Agent for Iterative and Fine-Grained Image Editing.}
RefineEdit-Agent is a training-free agent for multi-step editing: an LVLM parses instructions into subgoals, an LLM selects/parameterizes tools (ControlNet, GLIGEN, InstructPix2Pix, etc.), and an LVLM evaluates fidelity/preservation after each edit, feeding back for re-planning until a threshold or budget is met. ~\cite{avg250817435}

\paragraph{MIRA.}
MIRA is a lightweight VLM agent that decomposes complex editing instructions into atomic steps via an iterative perception-reasoning-action loop: it observes the current image and instruction, emits one atomic edit or STOP, executes it via an external editor, and repeats until termination. ~\cite{avg251121087}

\paragraph{JarvisEvo.}
JarvisEvo is a self-evolving photo-editing agent that jointly trains editor and evaluator policies via synergistic policy optimization on 170K Lightroom-style editing traces. It uses interleaved multimodal reasoning and self-generated reflection data to refine operations across 200+ tools. ~\cite{avg251123002}

\paragraph{Agentic Retoucher.}
Agentic Retoucher is a perception–reasoning–action agent that detects and corrects localized defects (hands, faces, text, geometry) in T2I outputs. A perception agent predicts distortion saliency masks, a reasoning agent diagnoses defect types, and an action agent selects mask- or instruction-guided inpainting; the loop repeats until artifacts are resolved, typically within 2–3 iterations. ~\cite{avg260102046}

\paragraph{EditRefiner.}
EditRefiner is a four-agent perception–reasoning–action–evaluation loop that detects and corrects local artifacts in text-guided image edits. A perception agent predicts artifact/failure saliency maps, a reasoning agent diagnoses flaw types, an action agent performs masked re-editing, and an evaluation agent scores quality and decides whether to continue. The loop converges in ~2.5 iterations. ~\cite{avg260507457}

\subsection{Planning and Tool Use for Image Creation}
Planning and tool-use systems make an image task executable through model routing, external knowledge, structured controls, toolpaths, and task reformulation. The distinction below is between choosing a grounded plan and constructing an operational workflow.

\subsubsection{Grounded Planning, Model Routing, and Structured Control}
Planning-oriented systems first make missing knowledge, spatial constraints, model capabilities, or control variables explicit. Some also revisit the artifact after rendering, whereas others remain pre-generation routing systems.

\paragraph{DiffusionAgent.}
 routes text-to-image requests to domain/style specialist models. An LLM parses the input (instruction, inspiration, or hypothesis), searches a subject–style Tree-of-Thought model index, consults an advantage database built from 10,000 prompts and human feedback, selects an expert model, and extends the prompt from model-specific examples. ~\cite{qin2024diffusionagent}

\paragraph{Divide and Conquer.}
CompAgent is an LLM-driven training-free framework for compositional T2I (text-to-image) generation. It decomposes prompts into objects/attributes with bounding boxes, selects among customization, layout-to-image, and localized editing tools, and uses multimodal inspection to repair incorrect attributes/relations. ~\cite{wang2024divide}

\paragraph{IA-T2I.}
 retrieves reference images from the web when T2I prompts contain uncertain, rare, or ambiguous knowledge. An active retrieval module decides whether to search, a hierarchical selector ranks results, and self-reflection scores outputs and triggers re-retrieval/generation if needed. ~\cite{li2025iat2i}

\paragraph{World-To-Image.}
 grounds T2I generation with web-retrieved knowledge for prompts involving novel/long-tail entities. An orchestrator diagnoses concept risk and selects among semantic decomposition, substitution, or image retrieval to condition OmniGen2. ~\cite{avg251004201}

\paragraph{GenAgent.}
 is an agentic multimodal model that iteratively reasons, generates, judges, and reflects over multiple-turn interactions, treating image generators as callable tools until the output satisfies the request. ~\cite{jiang2026genagent}

\paragraph{Mind-Brush.}
 is an agentic framework that turns T2I generation into a dynamic and knowledge-driven workflow, Simulating a human-like ’think-research-create’ paradigm: it detects knowledge/reasoning gaps from user intent, routes them to web/image search or CoT reasoning, accumulates evidence, and compiles a master prompt for final generation. ~\cite{avg260201756}

\paragraph{DiffGraph.}
 is an agent-driven framework that merges online expert models (checkpoints and LoRAs) for diverse T2I prompts. A graph construction agent registers expert descriptions as nodes and calibrates their capabilities on reference prompts as edge features; at inference, an expert selection agent retrieves relevant experts and a VGAE predicts merging coefficients from the activated subgraph, producing a merged model for rendering.~\cite{avg260320470}

\paragraph{Gen-Searcher.}
 is the first trained search-augmented image
generation agent. It performs multi-hop web search, image search, and browsing over multiple turns, then emits a grounded prompt with ordered reference images for a single generator call. ~\cite{avg260328767}

\paragraph{MetaPoint.}
 enables precise spatial control in UMMs via special tokens that reuse native 2D coordinate encodings: one token specifies a point, two define a box, and sequences encode layouts or edit regions; a VLM planner localizes objects and issues executable commands. ~\cite{zhou2026metapoint}

\paragraph{One Image is All You Need (WMGen-v1).}
WMGen-v1 generates long-tail training data for spatial perception from a single reference image: a Large
Vision-Language Model (LVLM) extracts scene semantics, camera properties, and spatial constraints; an LLM expands these under physical/commonsense guidance; a diffusion model renders synthetic variants. ~\cite{avg260620764}

\paragraph{RS-Gen.}
 is a training-free multi-stage agent for knowledge-intensive T2I generation/editing. An image router resolves coreferences against dialogue history; an intent analyzer routes to direct generation or structured questioning; a ReAct-based reasoning-and-search agent invokes web/image/geographic/VQA/reasoning tools with fallback retrieval; a generation agent executes a generate–verify–correct loop. ~\cite{avg260623221}

\paragraph{Qwen-Image-Agent.}
 bridges the context gap in T2I generation by constructing full generation context from underspecified user input via three-tier planning (information/content/generation-level) and four-source grounding (reasoning, search, memory, feedback). Its unified context-centric design enables multi-image and multi-turn generation while preventing context drift.~\cite{avg260626907}

\subsubsection{Workflow Construction and Task Reformulation}
These systems choose and sequence heterogeneous image operations. Their primary object of control is a toolpath, workflow graph, or reformulated editing task rather than a single renderer call.

\paragraph{CAISE}
 is a dataset and task for conversational image search and editing, where an agent generates executable commands (search, color/brightness/contrast/rotation/background removal) from dialogue and visual context. The dataset contains 1,611 dialogues and 6,173 commands; the baseline generator-extractor model achieves 46.43\% exact accuracy vs. 90.0\% for human experts. ~\cite{kim2022caise}

\paragraph{GenArtist.}
 is a unified image generation/editing system driven by an MLLM agent. It decomposes complex prompts, builds a planning tree of generation/editing/auxiliary nodes with visual verification after each generation, selects tools from a heterogeneous library, supplies missing positions via detectors/segmentation/pose-depth preprocessors, and backtracks to sibling tools on failure.~\cite{wang2024genartist}

\paragraph{RestoreAgent.}
 is an MLLM-driven agent that selects task order and models for mixed image degradations (denoising, deraining, dehazing, deblurring, JPEG removal, low-light enhancement). A vision encoder + Llama3 controller predicts the next tool from current image and history, with step-wise reassessment and rollback on failure. ~\cite{avg240718035}

\paragraph{AgenticIR.}
 is an LLM/VLM-driven system that restores images with multiple interacting degradations (noise, blur, compression...) by composing specialist models via a perception–scheduling–execution–reflection–rescheduling loop; failed operations trigger rollback and re-planning. A fine-tuned DepictQA provides degradation descriptions and success judgments; offline exploration enumerates tool sequences on held-out images and summarizes success patterns into retrievable documents for scheduling.~\cite{avg241017809}

\paragraph{HybridAgent.}
 is an agentic image restoration system with a FastAgent for direct prompts (routing tools in ~12\% of SlowAgent time) and a fine-tuned SlowAgent for ambiguous requests; a FeedbackAgent judges cleanness and can trigger further restoration. Tools are built via three-stage training (base model → LoRA single-task tools → mixed-degradation tool), enabling joint restoration to replace damaging long chains.~\cite{avg250310120}

\paragraph{CoSTA*.}
 is a cost-sensitive agent that plans and executes multi-turn image edits via LLM-driven subtask tree generation, tool subgraph construction over 24 tools, and A* search balancing quality and runtime with a tunable cost-quality parameter; a VLM validator checks each subtask and triggers retries or alternative paths on failure. ~\cite{gupta2025costa}

\paragraph{ComfyGPT}
 is a multi-agent system that generates executable ComfyUI workflows from natural-language task descriptions. ReformatAgent converts JSON workflows to link-centered diagrams, FlowAgent predicts diagrams (trained with SFT + GRPO with validity rewards), RefineAgent retrieves current node definitions from a database to replace invalid nodes, and ExecuteAgent submits the final workflow. ~\cite{huang2025comfygpt}

\paragraph{Q-Agent.}
 is a quality-driven image restoration (IR) agent that addresses multiple degradations via CoT-based perception and greedy IQA-guided restoration. A fine-tuned MLLM answers per-degradation Yes/No questions; at each step, all candidate operations are applied, ranked by an aggregate of five NR-IQA metrics, and the best is greedily selected until no improvement. ~\cite{avg250407148}

\paragraph{FaSTA*.}
 extends CoSTA* by mining reusable tool sequences (subroutines) from prior traces via LLM-based inductive reasoning and storing them with activation conditions in a rule table. It first attempts fast planning using these rules, with VLM checks validating each subroutine; only when no rule applies or a check fails does it fall back to localized A* search. ~\cite{gupta2025fasta}

\paragraph{Restore-R1.}
 is a lightweight reinforcement-learned agent that predicts restoration tool sequences in a single forward pass, eliminating runtime reflection/rollback. A CLIP+MLP policy is trained via PPO/GRPO using DeQA-Score as the reward, without ground-truth labels. ~\cite{avg251218599}

\paragraph{Derain-Agent.}
 is a plug-and-play post-deraining enhancer that refines initial derained outputs by predicting one of 16 predefined tool paths (denoising, deblurring, color correction) and pixel-wise strength maps from a ResNet34 feature extractor, then executing the path once. ~\cite{avg260311866}

\paragraph{TIR-Agent.}
is a trainable VLM agent that learns direct tool-calling policies for multi-degradation restoration via SFT + RL, replacing costly heuristic search. A Qwen3-VL-8B controller selects task and model per step from current image and history; EDP diversifies SFT trajectories, while MAR dynamically reweights FR/NR metrics to prevent reward hacking. ~\cite{avg260327742}

\paragraph{IMAGAgent.}
 is a plan–execute–reflect agent for multi-turn image editing that prevents error accumulation and semantic drift. A VLM planner decomposes instructions into atomic sub-tasks; a controller dynamically composes segmentation, detection, retrieval, and editing tools from the current image and history; multiple VLM experts critique intermediate outputs; and an aggregator converts feedback into corrective actions, with up to three retries per turn. ~\cite{avg260329602}

\paragraph{Adaptive Task Reformulation.}
ATR improves image editing by reformulating ambiguous tasks without modifying the backbone. A profiler extracts target and scene context; a router selects direct/rewritten editing, spatial decoupling, or localized workspace; a planner executes sequentially with state, feedback, and fallback. ~\cite{avg260415917}

\paragraph{From Plans to Pixels.}
 is an experiential learning framework for long-horizon advertisement editing. A checklist-guided planner decomposes abstract instructions into subtasks; a reward-trained orchestrator selects tools/regions from a heterogeneous editor library; a VLM judge provides outcome rewards, infeasible subtasks are pruned, and a distilled verifier re-ranks intermediate candidates. ~\cite{avg260515181}

\paragraph{OPERA}
 is an end-to-end framework for multi-tool cooperative image restoration. A GRPO-trained planning agent emits complete tool compositions directly from the input, receiving final restoration quality as reward; 16 restoration tools are jointly fine-tuned through generated chains, enabling cooperative behavior. ~\cite{avg260522104}

\paragraph{GenClaw.}
 is a code-driven agentic image generation framework that decouples creation into Conceptualize–Sketch–Color stages: it retrieves contextual knowledge via search/reasoning, constructs executable visual code (SVG/HTML/Canvas/Three.js) as an intermediate layout, then renders photorealism via an image generator. ~\cite{ye2026genclaw}

\paragraph{DiTTo.}
\textbf{Problem and setting.} DiTTo addresses order-aware restoration when the pool of restoration experts expands and generating real-expert trajectories for every possible order becomes prohibitive~\cite{avg260530915}. \textbf{Method and mechanism.} A simulator predicts the one-step effect of each restoration action and an all-in-one IQA model ranks candidate actions, reducing synthetic trajectory construction to a linear number of simulator calls. A vision-language agent is first trained on these trajectories and then aligned to a small real-expert trajectory set by decomposed preference optimization over degradation perception, order selection, and valid tool calls. When a new expert is added, the reusable simulator, evaluator, and SFT checkpoint are retained while only the alignment stage is updated. \textbf{Evaluation and results.} MiO-100 contains 2,400 test images with two to five concurrent degradations. With the common JarvisIR tool pool, DiTTo obtains MUSIQ/MANIQA/CLIP-IQA+/NIQE of 69.40/0.6457/0.7762/5.475 for two degradations and 69.34/0.6270/0.7690/5.224 for five; the extended pool reaches 71.76/0.7127/0.8393/5.208 and 72.27/0.7241/0.8509/5.188. The reported cost of adapting to a new restoration expert is about 60 hours on two B200 GPUs, compared with about 907 hours for JarvisIR under the matched setup. \textbf{Evidence boundary.} Sequential current-state tool calls supply L3 correction evidence, and the demonstrated reuse and update of the alignment policy for a new expert provide L5 evidence for cross-task mechanism improvement. This evidence is offline adaptation rather than unattended online self-evolution; simulator-to-expert mismatch remains material, and reported quality evaluation is entirely no-reference.

\paragraph{IEA.}
\textbf{Problem and setting.} IEA is an amateur-friendly conversational editor for global photo retouching tasks such as brightness, exposure, contrast, temperature, saturation, highlights, and shadows. The artifact is an edited photograph accompanied by an interpretable sequence of parameterized operations and, in the refinement task, a later user instruction~\cite{avg260608016}. \textbf{Method and mechanism.} A Qwen2.5-VL-7B policy operates 16 tools in a simulated editor. Stage 1 uses distilled expert programs from GIER for supervised alignment; Stage 2 applies GRPO with likeness and usefulness rewards; Stage 3 adds synthetic Image-Edit, Image-Summary, and Image-Refine data so the policy can summarize edit history and adjust parameters in response to user feedback. The model emits explicit tool calls rather than directly synthesizing pixels, and prior operations are supplied as the refinement context. \textbf{Evaluation and results.} The real GIER-derived training set contains about 2.4k image pairs after filtering, while synthetic data expand the mixed supervision to about 202k examples. On 2,568 Image-Edit and 241 Image-Summary test samples, IEA-Stage-3 reaches pixel distance $0.103$, likeness reward $0.149$, usefulness reward $0.402$, ROUGE-L $0.258$, and alignment reward $7.387$, improving over Stage 1 and Stage 2. In a 50-sample user study with 56 participants, IEA obtains the best rank among tool-calling systems for instruction following and image quality, with ranks $4.64$ and $3.69$. \textbf{Evidence boundary.} The history-aware refinement action can respond to a user's assessment of a previous edit, supporting L3 human-in-the-loop feedback adaptation. The runtime uses a fixed parameterized tool inventory and a simulated editor, and it does not automatically diagnose the rendered image, roll back branches, or update policy from completed user sessions. The evidence therefore does not establish L4 long-horizon autonomy or L5 cross-task self-improvement.

\paragraph{CanvasAgent.}
\textbf{Problem and setting.} CanvasAgent addresses complex image creation and editing requests that combine generation, grounding, segmentation, extraction, compositing, OCR, geometric transformation, and enhancement across several input and intermediate assets~\cite{avg260705465}. \textbf{Method and mechanism.} CanvasCraft supplies 140,000 executable supervised trajectories and 10,000 reinforcement-learning tasks over 11 visual tools. CanvasAgent first learns valid reasoning--action formats by supervised fine-tuning and then uses GRPO with outcome, process, rule, and efficiency rewards. During execution it maintains explicit asset identifiers, observes intermediate outputs, and can revise the plan, switch tools, verify with OCR or grounding, or return to an earlier image state. \textbf{Evaluation and results.} On the 250-task CanvasCraft evaluation split, SFT plus RL reaches 0.821 overall reward, 0.869 alignment, 0.849 trajectory quality, and 0.785 rule-based score, compared with 0.557, 0.613, 0.576, and 0.467 for SFT alone. In a human evaluation over 12 tasks, it scores 3.97 for task alignment, 3.90 for key details, and 4.06 for aesthetics; Qwen3-VL-32B scores 2.93, 2.91, and 2.83. \textbf{Evidence boundary.} Multi-asset state, observation-dependent tool choice, verification, and rollback support L4 long-horizon control. The tool inventory remains fixed, RL requires real tool execution and an external judge, and the human study is small. Training-time policy optimization is completed before deployment; no reusable experience is acquired from completed user tasks, so the evidence does not meet L5.

Workflow construction expands the action space available for image creation, while its reliability depends on valid tool interfaces, state carried between operations, and evidence that a failed path can be revised. Several studies learn policies offline; that alone does not demonstrate cross-task continual improvement at deployment.

\subsection{Stateful Image Generation and Editing}
Stateful image systems retain information needed to preserve identity, user intent, scene structure, and earlier decisions. The literature separates turn-level user and context management from explicit scene, narrative, and trajectory representations.

\subsubsection{Multi-Turn User, Subject, and Context Management}
These systems retain information that is needed across turns, including user preferences, subject identity, prior operations, context, and candidate quality. The persistent representation makes a later action dependent on more than the current prompt.

\paragraph{AutoStudio.}
\textbf{Problem and setting.} AutoStudio studies multi-turn interactive image generation in which users may introduce, customize, refer back to, or edit several subjects over successive turns; the central requirements are cross-turn identity consistency, correct interactions, and support for on-the-fly instructions~\cite{cheng2024autostudio}. \textbf{Method and mechanism.} The training-free framework stores subject and component features with persistent identifiers in a subject database. A subject manager converts the dialogue history and current request into subject-aware captions, a layout generator predicts component-level bounding boxes, and a supervisor critiques spatial relations and returns suggestions until the layout is refined. A diffusion-based drawer then combines the layout with retrieved subject features through P-UNet and subject-initialized generation. \textbf{Evaluation and results.} Evaluation on CMIGBench uses average Fr\'echet Inception Distance, character--character similarity, and text--image similarity, supplemented by human assessment. The paper reports improvements over the previous best results of 13.65\% in average FID and 2.83\% in average character--character similarity, and its ablation shows a marked reduction when the supervisor is removed. \textbf{Evidence boundary.} The supervisor's critique changes the subsequent layout, and earlier generated subjects are carried into later turns, which supports L3 feedback adaptation over structured layout and interaction state. The automatic supervisor does not diagnose a completed image and select a localized repair action; image-level failures and the quality of human-directed later turns therefore remain outside the demonstrated automatic loop.

\paragraph{Preference Adaptive and Sequential Text-to-Image Generation (PASTA).}
\textbf{Problem and setting.} PASTA studies collaborative text-to-image creation where a user's initial prompt underspecifies an evolving visual preference~\cite{nabati2025pasta}. \textbf{Method and mechanism.} A reinforcement-learning agent uses the history of user selections and a learned preference model to present adaptive prompt-expansion slates over multiple turns. The work collects human sequential-preference data and uses it with simulated trajectories to train the policy. \textbf{Evaluation and results.} The study collects over 7,000 five-step human-rater sequences from about 100 raters and augments them with more than 30,000 simulated rollouts. The learned user model reaches about 70\% accuracy on its human-rated metrics. In direct final-turn comparison, PASTA reports an 85\% relative improvement over the Gemini~1.5 Flash baseline, and it reports 20\% relative improvement when transferred from SDXL to Flux without retraining. \textbf{Evidence boundary.} Sequential user choices and preference estimates alter later prompt expansions, supporting L3 feedback adaptation. Preference state is learned offline and the rater study notes potential group and timing bias; the work does not demonstrate a deployed controller that continually updates from later user sessions.

\paragraph{StoryState.}
\textbf{Problem and setting.} StoryState addresses illustrated storybooks in which character identity, world settings, and page composition must remain consistent while a local or global edit changes only the intended pages~\cite{sarkar2026storystate}. \textbf{Method and mechanism.} A Planner and State Manager externalize the book as $S=(C,W,\{S_i\})$, comprising a character sheet, global world state, and page-level scene states with references to text and image assets. A Prompt Writer maps this persistent state to global and page prompts. Local edits update one $S_i$; identity edits update $C$ and propagate only to affected pages. After regeneration, a multimodal Consistency Critic checks the new page against the state and neighboring pages and proposes a minimal state correction, which is applied after acceptance. \textbf{Evaluation and results.} On 192 ten-page storybooks, StoryState reaches 0.83 adjacent-page consistency, compared with 0.78 for 1Prompt1Story, while reducing pages changed per edit from 4.5 to 1.6, interaction turns from 4.3 to 3.1, and edit time from 96 to 74 seconds. Gemini Storybook reaches 0.89 consistency but provides no iterative editing. In a 100-participant study, StoryState receives the highest preference for consistency (36\%) and control (48\%). \textbf{Evidence boundary.} Explicit persistent state and selective propagation across a long image sequence support L4 stateful control. Correction remains partly human governed because critic proposals are applied after acceptance, and removing the critic changes consistency only from 0.83 to 0.81 in the reported ablation. The framework does not learn a reusable policy across books and inherits identity and rendering limitations from its text-to-image backend.

\paragraph{Agent Banana.}
\textbf{Problem and setting.} Agent Banana targets professional image-editing sessions in which vague or composite instructions must be resolved into localized operations while preserving non-target content, object state across turns, and native 4K detail~\cite{ye2026agentbanana}. \textbf{Method and mechanism.} A hierarchical planner--executor architecture decomposes requests into atomic add, remove, replace, adjust, and undo operations. Context Folding compresses long interaction histories into structured state, Image Layer Decomposition isolates a high-resolution object layer for local editing and lossless fusion, and a visual evaluator can trigger retry, rollback, or replanning. The state graph retains earlier images, making undo and branch-specific repair executable actions rather than new global generations. \textbf{Evaluation and results.} HDD-Bench contains 96 curated three-turn sessions with verifiable object-state transitions and native 11.8-megapixel images. Agent Banana reports image consistency of 0.871, masked SSIM of 0.84, and masked LPIPS of 0.12; Nano Banana Pro obtains 0.861, 0.72, and 0.14 under the same table. Its instruction-following score is 0.849, below Nano Banana Pro's 0.911, while the reported single-turn editing scores remain competitive. \textbf{Evidence boundary.} Intermediate verification, persistent state, rollback, replanning, and evaluator-controlled continuation establish L3 feedback adaptation and implement several L4 mechanisms. We retain L3 because the controlled benchmark contains only three turns and does not evaluate authority escalation or extended autonomous stopping and recovery. Performance also depends on GPT-5-mini and proprietary editing tools, and the graph-based state estimator can miss visual changes outside its object attributes.

\paragraph{Generation Navigator.}
\textbf{Problem and setting.} Generation Navigator addresses the choice between preserving a satisfactory image, editing a partially correct one, and discarding a structurally flawed result during multi-turn image creation~\cite{avg260517969}. \textbf{Method and mechanism.} A learned multimodal navigator observes the original request, prior actions, generated images, reviewer critiques, and scalar scores. It selects \textsc{Stop}, \textsc{Refine}, or \textsc{Regenerate}, together with a revised prompt; the generator executes the action and the reviewer supplies the next observation. Training combines 103,000 constructed trajectories, supervised fine-tuning, and PRE-GRPO, whose trajectory reward accounts for peak quality, non-regression, turn efficiency, and valid action format. The reported deployment budget is three turns, with the highest-scored image retained if the budget expires. \textbf{Evaluation and results.} With a Qwen3-VL-8B navigator, FLUX.2-Klein-9B generator, and Doubao-Seed1.5 reviewer, the system obtains 79.06% reasoning accuracy and 94.83 image quality on T2I-ReasonBench, compared with 57.70% and 82.82 for the generator alone. It reaches 0.90 overall on WISE, compared with 0.61 for the same FLUX generator, and 0.8843 on GenEval, compared with 0.8481. Transfer experiments raise Qwen-Image reasoning performance by 21.66 points, and replacing the inference-time reviewer produces a 3.03-point range. The full chain-of-thought configuration records 77.4% accuracy at 52.5 seconds per sample; a no-chain-of-thought variant records 76.0% at 21.7 seconds. \textbf{Evidence boundary.} Current images and reviewer feedback directly determine later edit, regeneration, and termination decisions, providing L3 feedback adaptation. The policy is trained before deployment and its trajectory is capped at three turns; the paper does not establish persistence beyond the request, rollback to an earlier state other than incumbent selection, or continuous improvement from completed tasks. Its correctness therefore remains sensitive to reviewer diagnostics and the available generator/editor.

Multi-turn state is useful when it selectively constrains later generation or editing actions. The remaining question is whether state is durable, auditable, and recoverable rather than merely appended conversation context; most reported systems retain it only within one task.

\subsubsection{Scene, Narrative, and Trajectory State Management}
For multi-object edits and visual narratives, systems externalize scene structure, causal relations, candidate branches, or operation histories. This turns an image sequence into an addressable state space for planning and recovery.

\paragraph{I2E.}
\textbf{Problem and setting.} I2E targets compositional image editing with multiple instances, spatial relations, occlusions, and physical dependencies. Its artifact is an editable scene state rather than an undifferentiated pixel canvas, so actions such as moving, removing, resizing, inserting, or recoloring one object can preserve other instances~\cite{avg260103741}. \textbf{Method and mechanism.} A Decomposer segments relevant objects, reconstructs occluded foreground and background regions, and builds a physically ordered layer environment with a DAG of spatial constraints. A VLA Editor converts a language request and current scene state into atomic parameterized actions. Physics-aware reasoning checks relations such as gravity and support; each action updates the environment, and user feedback or self-critique is appended as a corrective action without resetting the scene. The resulting loop executes, evaluates, and revises multi-action edits over the maintained state. \textbf{Evaluation and results.} I2E-Bench contains 200 images with 5--10 instructions each. On this benchmark, I2E reports LPIPS-U $0.0754$, DINO $0.9821$, spatial accuracy $0.6923$, constraint satisfaction $0.8700$, physical consistency $0.9210$, instruction compliance $0.8645$, and multi-step score $0.8074$. On MagicBrush it reaches $0.0446/0.9581/0.7120/1.0000/0.8668/0.7924$ for the reported fidelity, consistency, spatial, constraint, physical, and instruction measures; on EmuEdit it reaches $0.0565/0.9404/0.7183/1.0000/0.8952/0.8107$. A blind study with 30 participants gives I2E an average score of $7.42$ and mean rank $1.79$. \textbf{Evidence boundary.} The explicit scene state, dependency-aware action sequencing, corrective action accumulation, and multi-instance demonstrations support L4 stateful control. Comparative experiments restrict all methods to one refinement round, so the quantitative evidence for repeated autonomous recovery is narrower than the full runtime design. Segmentation and inpainting errors propagate through later actions, physical simulation is limited for soft bodies and fluids, and no cross-task policy or skill update is reported.

\paragraph{MSRAMIE.}
\textbf{Problem and setting.} MSRAMIE targets image editing with long, multi-instruction requests whose subgoals compete for attention and whose repeated execution can damage identity and perceptual quality. It treats intermediate edited images as candidate states in a constrained search space~\cite{avg260316967}. \textbf{Method and mechanism.} An MLLM Instructor and a plug-in editing Actor build a Tree-of-States. Each state stores its input image, output image, editing thought, and evaluation scores, while root-recall links retain the original instruction and image. A Graph-of-References retrieves nearby states to avoid duplicate exploration. The Instructor evaluates instruction following and identity preservation, generates a prompt focused on unmet requirements, and a scheduler chooses expansion, traceback, resampling, aggregation, or termination when a completion state, budget, or topology limit is reached. \textbf{Evaluation and results.} Complex-Edit supplies 531 real images across seven instruction-complexity levels. With Qwen-Image-Edit, MSRAMIE raises overall VQAScore from $0.7623$ to $0.8819$, complete instruction following from $0.3497$ to $0.6008$, and CLIP-I from $0.9294$ to $0.9335$, while overall FID changes from $116.51$ to $116.69$. At complexity seven, VQAScore rises from $0.7172$ to $0.8780$ and complete instruction following from $0.1676$ to $0.4708$. Similar improvements are reported with Flux-Kontext and Flux-Klein. In the ablation, Tree-of-States reaches VQAScore $0.8666$, and adding the Graph-of-References raises it to $0.8819$. \textbf{Evidence boundary.} Persistent candidate states, traceback, cross-branch references, metric-conditioned scheduling, and budget-aware termination provide L4 evidence. The process is training-free and stores state only within a request; repeated encoding can reduce perceptual quality, evaluations use an MLLM-derived checklist, and the number of states grows approximately linearly with instruction complexity. No cross-task learning is demonstrated.

\paragraph{LogiStory.}
\textbf{Problem and setting.} LogiStory targets multi-image narratives whose individual panels may be visually plausible while actions, object states, emotions, and causal consequences remain unclear or contradictory across the sequence~\cite{avg260328082}. \textbf{Method and mechanism.} SceneCrafter, LogicMiner, and ShotPlanner agents construct entity definitions, action--state events, and panel specifications. A Local Causal Monitor maintains a text-form memory of accepted panels and scores each new image against the accumulated reading path. A Global Causal Verifier maintains a story-level causal graph and dynamic instance states. Depending on two thresholds, the system accepts the image, regenerates it, or invokes visual editing with verifier instructions before updating memory. \textbf{Evaluation and results.} On the 60-story LogicTale benchmark, LogiStory obtains 4.23 for instance consistency, 4.45 for narrative causality, and 0.8267 for story readability. GPT-4o plus GPT-Image-1 obtains 4.67, 3.96, and 0.7622, respectively. A 30-participant evaluation gives LogiStory 4.25 for instance consistency, 4.26 for causality, and 3.83 for readability. Using both local and global verification raises readability from 0.6527 with neither module to 0.8267. \textbf{Evidence boundary.} Persistent causal memory, dynamic instance state, conditional regeneration or editing, and per-panel acceptance provide L4 evidence for stateful sequence control. The framework does not retain experience across stories, and its modular pipeline adds computation and depends on proprietary language, generation, and editing models. High-entity stories and subtle emotional transitions remain documented failure cases.

Explicit structured state permits selective propagation, traceback, and localized correction. Its practical value depends on the accuracy of segmentation, scene reconstruction, and verifier judgments that update the state.

\subsection{Multi-Agent Systems for Image Generation and Editing}
Multi-agent systems allocate visual creation work among specialized roles. Their value depends on how role outputs, shared state, and evaluator feedback determine later actions, rather than on the number of participating agents.

\subsubsection{Collaborative Generation, Composition, and Domain Design}
This group distributes generation, analysis, design, or evaluation work across specialized roles. It includes contextual comparisons that clarify the difference between model-level cooperation and task-level visual creation control.

\paragraph{Message Passing Multi-Agent GANs.}
\textbf{Problem and setting.} This early work studies unsupervised image generation with multiple generators, motivated by the limited communication and mode specialization available in a conventional single-generator GAN~\cite{avg161201294}. \textbf{Method and mechanism.} Two generators share a discriminator and exchange learned messages through a common message generator; conditioned message passing combines the other generator's output, its latent input, and the previous message, while competing and conceding objectives encourage complementary behavior. \textbf{Evaluation and results.} Experiments use SVHN representations and CelebA message clusters. The concatenated representation of the different-noise conditioned model reduces SVHN classification error from the reported DCGAN value of 22.48\% to 15.2\%, although the Improved GAN reference remains lower at $8.11\pm1.3$\%; interpolation and t-SNE analyses indicate that the two generators capture different visual traits. \textbf{Evidence boundary.} The agents are generator modules trained in a fixed architecture. There is no task-level planning, tool choice, artifact diagnosis, or recovery policy, so the evidence supports the L1 baseline in this survey rather than a modern closed-loop visual agent.

\paragraph{Anywhere.}
\textbf{Problem and setting.} Anywhere studies foreground-conditioned image generation, where a supplied object must be placed in a new background without damaging its silhouette, creating a semantically incompatible scene, suppressing output diversity, or ignoring optional textual instructions~\cite{xie2024anywhere}. \textbf{Method and mechanism.} A foreground analyzer and prompt-generation agents describe the object and select a compatible scene prompt. A template generator produces the background, after which a template-repainting agent detects mask disagreement around the protected object and repairs the affected region. A high-resolution enhancer composites and harmonizes the foreground. Finally, a vision-language quality evaluator inspects the completed image, returns textual feedback to the prompt module, and requests another attempt when quality is inadequate; the loop is capped at three rounds. \textbf{Evaluation and results.} The authors construct a 3,000-image foreground test set from LAION and MS COCO. In the text-free setting, Anywhere reports FID 37.007 and PickScore 0.418, compared with 38.797 and 0.133 for HD-Painter and 41.047 and 0.261 for BrushNet. In the text-guided setting, its FID is 36.450 and CLIP similarity is 18.357. A 10-participant study over 100 foreground cases reports a diversity score of 2.57 and a bad-case rate of 0.19, compared with 1.93 and 0.34 for the strongest corresponding baseline values. \textbf{Evidence boundary.} The final-image evaluator can alter the next prompt and trigger regeneration, which supports L3 artifact-conditioned correction. The maximum of three rounds, approximately doubled GPU time, dependence on several external models, and failure cases involving transparent foregrounds limit the evidence for long-horizon recovery or persistent autonomy.

\paragraph{Marmot.}
\textbf{Problem and setting.} Marmot targets object-counting, attribute-binding, and spatial-relation errors in complex text-to-image outputs, where serial whole-image editing can omit objects or accumulate background distortion~\cite{sun2025marmot}. \textbf{Method and mechanism.} An Object-Aware Agent decomposes the requested scene into object-centric correction tasks. For each object or object pair, an Object Correction System combines decision, execution, and verification over a segmentation mask or pair of bounding boxes; failed verification restores the original region and retries with another seed. A Pixel-Domain Stitching Smoother merges independently corrected regions through mask-guided latent optimization. \textbf{Evaluation and results.} On T2I-CompBench, adding Marmot to SDXL changes the color, shape, texture, and spatial scores from 0.5673, 0.4958, 0.5865, and 0.2035 to 0.6849, 0.5527, 0.6412, and 0.3347. On GenEval, the SDXL overall score rises from 0.55 to 0.66; removing verification reduces the T2I-CompBench color score from 0.6849 to 0.6339. \textbf{Evidence boundary.} Object-level verifier decisions trigger retries and affect the assembled image, supporting L3 feedback adaptation. Correction depends on successful object localization and available editors, and missing target objects cannot be repaired reliably in the reported layout-to-image experiments. The process uses a fixed decomposition and correction repertoire without cross-task learning.

\paragraph{MCCD.}
\textbf{Problem and setting.} MCCD targets complex text-to-image prompts containing multiple objects, attributes, actions, depth order, and spatial relations that a single diffusion condition often binds incorrectly~\cite{avg250502648}. \textbf{Method and mechanism.} Its Multi-Agent Collaboration-based scene Parsing module contains a conductor, evaluator, and specialist agents for objects, background, action relations, spatial relations, layout, and aesthetics. The conductor dynamically chooses the next specialist during forward reasoning. The evaluator checks the accumulated structured prompt set; when it identifies an erroneous component, backward feedback localizes the responsible agent, updates its output, and restarts forward reasoning. A Hierarchical Compositional Diffusion module then maps object and background prompts to depth-aware Gaussian regions, enhances their latent contributions, and smooths box boundaries. \textbf{Evaluation and results.} T2I-CompBench evaluation shows improvements across four Stable Diffusion backbones. For SDXL-base, adding MCCD raises color binding from 0.5744 to 0.6278, texture binding from 0.4907 to 0.5647, spatial-relation accuracy from 0.1971 to 0.2350, and the complex score from 0.3130 to 0.3348. Component ablations and examples with alternative multimodal language models further test the parsing and diffusion modules. \textbf{Evidence boundary.} Intermediate evaluator evidence can change which parsing agent is revisited and what structured condition reaches the renderer, supporting L3 feedback adaptation. The correction loop is concentrated in prompt and layout construction; the paper does not show the generated image being repeatedly diagnosed and repaired. Its quantitative evidence is limited to T2I-CompBench, and the method inherits the spatial and rendering capacity of the selected diffusion backbone.

\paragraph{T2I-Copilot.}
\textbf{Problem and setting.} T2I-Copilot addresses ambiguous or underspecified text-to-image requests and the repeated manual prompt editing needed when a generated image does not match user intent~\cite{chen2025t2icopilot}. \textbf{Method and mechanism.} The training-free system uses an Input Interpreter to extract entities, attributes, relations, and clarification questions; a Generation Engine to select and execute a generator or localized editing route; and a Quality Evaluator to score aesthetic and alignment criteria. Scores below a threshold produce structured suggestions that condition model reselection or regeneration, with optional user feedback and bounded retries. \textbf{Evaluation and results.} Experiments cover 200 DrawBench prompts and 1,600 GenAI-Bench prompts. With automatic LLM clarification, T2I-Copilot reaches a GenAI-Bench VQAScore of 0.813, compared with 0.745 for its FLUX.1-dev backbone and 0.588 for GenArtist; human clarification raises the score to 0.839. Removing the evaluator reduces the automatic system's score from 0.813 to 0.805, and a 2,442-vote user study reports average win rates of 94.5\% for alignment and 77.7\% for aesthetics. \textbf{Evidence boundary.} Image scores and diagnostic suggestions change later generation actions, supporting L3 feedback adaptation. The stopping rule is a fixed score threshold and retry cap, the evaluator relies on an MLLM's judgments, and the system does not retain corrections across independent requests.

\paragraph{From Image Generation to Infrastructure Design.}
\textbf{Problem and setting.} This work treats bicycle-lane visualization as a constrained street-image editing problem: a proposed intervention must follow roadway geometry and a specified design scenario while preserving vehicles, curbs, lighting, markings, and other non-target content~\cite{avg250905469}. \textbf{Method and mechanism.} A locator agent identifies whether a street view is suitable and returns curb, lane, and insertion-region polylines. A prompt agent combines this state with one of eight bicycle-lane specifications. The generation agent first paints a green intervention region and then converts it into a lane treatment with GPT-Image-1. For each scenario it produces ten candidates. A fine-tuned YOLO segmentation model isolates lane regions for CLIP reranking, and GPT-o3 checks the top three candidates against hard design constraints. If none passes, the system regenerates a new pool; experts can also correct descriptions, prompts, highlight regions, and upstream decisions. \textbf{Evaluation and results.} The study retrieves 600 Google Street View images from 150 Washington, D.C. sites and manually retains one view per site. The segmentation model, trained from 1,060 annotated real and generated images, reaches validation mAP@50 of 91.9\%. On a held-out street-view test set, evaluator selection accuracy ranges from 95.5\% to 97.0\% across the eight scenarios. Qualitative ablations associate removal of the locator with lane misplacement, removal of prompt optimization with boundary and buffer errors, and removal of the highlight stage with width drift and off-target edits. \textbf{Evidence boundary.} Candidate compliance changes selection and can trigger regeneration, supporting L3 correction. The system remains a domain-specific, staged workflow with substantial human authority; upstream revision is often expert initiated, the reported evaluator result does not establish fully autonomous design validity, and GPT-based verification adds roughly 60--90 seconds per scenario.

\paragraph{GenPilot.}
\textbf{Problem and setting.} GenPilot addresses semantic gaps in long or compositional prompts, including missing details, incorrect attributes, relations, counts, and negative constraints, across different text-to-image backbones~\cite{avg251007217}. \textbf{Method and mechanism.} An error-analysis stage decomposes the prompt into meta-sentences and combines VQA-based inspection with caption comparison to localize mismatches. A test-time optimization stage generates candidate rewrites, scores them with a multimodal model, clusters candidates, and uses a memory module to guide the next search round. The selected prompt is then sent to the chosen generator, with early stopping available. \textbf{Evaluation and results.} On the challenging DPG-bench subset, the reported average score improves from 72.04 to 74.08 for DALL-E 3, from 68.16 to 73.32 for FLUX.1 schnell, and from 53.16 to 62.12 for Stable Diffusion v1.4. On GenEval, FLUX.1 schnell increases from 65.82\% to 69.60\%, while the ablation without memory reaches 66.05\%. \textbf{Evidence boundary.} Error lists, candidate scores, and memory feedback alter later prompt candidates, supporting L3 feedback adaptation. The search operates in discrete text space and uses an MLLM scorer whose judgments can be unreliable; the released error patterns and model evaluations do not demonstrate persistent cross-task skill transfer, so the memory is evidence for within-task state rather than L5 continual improvement.

\paragraph{Collaborative Text-to-Image Generation.}
\textbf{Problem and setting.} This work studies multimodal text-to-image generation across architecture, portraiture, and landscape settings, where a single model may not preserve domain-specific vocabulary, visual detail, and text--image semantic alignment simultaneously~\cite{avg251010633}. \textbf{Method and mechanism.} It couples a text-enhancement subsystem with an image-generation subsystem, each containing domain-specialized agents and a multimodal fusion module. Proximal Policy Optimization trains the agents with a composite reward over semantic similarity, linguistic and visual quality, and diversity; contrastive learning and bidirectional attention connect the textual and visual representations. \textbf{Evaluation and results.} The study reports six experimental settings and evaluates linguistic output, visual quality, and cross-modal consistency. Transformer fusion attains the largest reported composite score, 0.521, while multimodal ensemble variants record consistency values between 0.444 and 0.481. The authors also report a 1,614\% increase in generated word count together with a 69.7\% reduction in ROUGE-1, exposing a tension between elaboration and lexical overlap rather than an unambiguous improvement. \textbf{Evidence boundary.} The contribution is a learned L2 specialization and fusion architecture. PPO modifies the agent parameters during offline training, while the paper does not show a deployed image observation that changes a subsequent rendering action, retained experience across user requests, or a benchmarked online self-improvement loop. Its evaluation is less comparable with current image-agent benchmarks because the reported composite metrics combine several generated modalities and the paper identifies instability in Transformer fusion.

\paragraph{Multi Agents Semantic Emotion Aligned Music to Image Generation.}
\textbf{Problem and setting.} MESA-MIG studies music-conditioned image generation in which visual content should reflect both musical semantics and continuous valence--arousal emotion, rather than relying only on low-level audio similarity~\cite{avg251223320}. \textbf{Method and mechanism.} An audio encoder and BART captioner extract instrumentation, rhythm, genre, and other semantic cues, while a parallel regression head predicts valence and arousal. Scene, verb, style, color, and composition agents transform these representations into visual attributes. A rule-based validator flags redundant, malformed, or affectively inconsistent attributes, and a prompt agent assembles validated alternatives for WanX 2.1-Turbo. The generated image is scored afterward by aesthetic, semantic, and cross-modal emotion models. \textbf{Evaluation and results.} The paper reports generation evaluation on 400 matched music--image pairs drawn from non-overlapping DEAM and EMOTIC partitions. MESA-MIG obtains an aesthetic score of 6.24/10 and semantic similarity of 0.895, compared with 5.88 and 0.870 for IMEMNet, 4.11 and 0.864 for SonicDiffusion, and 3.75 and 0.817 for Sound2Scene. Its music valence--arousal head also reduces valence RMSE from 0.238 to 0.150 relative to Music2Emotion. \textbf{Evidence boundary.} The specialists and validator form an interpretable L2 orchestration pipeline, but validation occurs before rendering and the measured output does not alter a subsequent generation action. The study depends on learned affect estimators, a curated cross-modal pairing procedure, and relatively small evaluation sets; it reports no human study of whether the generated image conveys the intended musical emotion and no deployed artifact-repair loop.

\paragraph{M3.}
\textbf{Problem and setting.} M3 addresses failures of object count, attribute binding, spatial relation, typography, and knowledge-dependent content in an initial text-to-image result or a user-supplied image~\cite{avg260206166}. \textbf{Method and mechanism.} Its training-free loop turns the request into a visual checklist, checks each constraint against the current image, converts a failed item into an actionable edit instruction, and invokes an instruction-guided editor. A verifier compares the edited candidate with the previous best image against the original request; rejected edits are retried up to two times, and an unresolved checklist item can be skipped to prevent an infinite loop. The full version uses Qwen3-VL-8B-Thinking for planning, checking, refinement, and verification with Qwen-Image-Edit as editor. M3-Hybrid replaces selected roles with spaCy rules and CLIP so that MindOmni can serve as checker and editor. \textbf{Evaluation and results.} On GenEval, M3-Hybrid raises MindOmni's overall score from 0.77 to 0.80, including count from 0.59 to 0.64 and attribute binding from 0.60 to 0.69. On the authors' harder GenEval split, M3 with a VLM planner raises Qwen-Image overall from 0.81 to 0.87 and position from 0.26 to 0.54. On OneIG-EN, M3 (VLM) reports 0.532 overall, with the rule-based configuration at 0.531. The ablations report complementary planner behavior and performance that generally saturates after three to five refinement iterations. \textbf{Evidence boundary.} Checklist failures, candidate comparison, and bounded retries change the next image-editing action, supporting L3 artifact-conditioned correction. The system does not retain a task memory or update its planner from completed work; its ability to diagnose constraints depends on the VLM or rule checker, while unresolved items may be skipped. The main reported quantitative results assess final compositional compliance rather than the reliability of failure localization or recovery across longer trajectories.

\paragraph{coDrawAgents.}
\textbf{Problem and setting.} coDrawAgents targets compositional prompts whose object counts, attributes, spatial relations, occlusions, and physical interactions are difficult to solve through one global layout or a single image sample~\cite{avg260312829}. \textbf{Method and mechanism.} An interpreter chooses direct layout-free generation for simple requests or decomposes a complex prompt into priority-ordered object groups. In layout-aware mode, a planner observes the partial image, previous layouts, and grounded existing objects before proposing the next placement. A checker evaluates local size and boundary conditions as well as global relations and plausibility, refines the proposal, and passes it to a painter that updates the canvas. The resulting partial image becomes evidence for the next priority group. \textbf{Evaluation and results.} On GenEval, coDrawAgents reports 0.94 overall, including 0.94 for counting and 0.95 for position, compared with 0.84 overall and 0.75 for position for GPT-Image-1 High. On DPG-Bench it obtains 85.17 overall, compared with 84.08 for SD3-Medium and 83.57 for OmniGen2. Across 1,074 DPG-Bench images, the interpreter, planner, checker, and painter are called 1.00, 1.52, 1.62, and 1.95 times per generation on average. \textbf{Evidence boundary.} The planner and checker use the evolving canvas to alter later layouts, which is direct L3 artifact-conditioned planning. The number of rounds is tied to semantic priority groups rather than learned stopping, and the paper reports no persistent memory, rollback across completed groups, or cross-task adaptation. Performance remains bounded by the integrated FLUX and 3DIS painters, while repeated multimodal calls add latency relative to direct generation.

\paragraph{AI-Gram.}
\textbf{Problem and setting.} AI-Gram is a continuously operating, image-centric social environment used to study how autonomous visual agents form reply chains, social ties, and collective visual patterns without human-authored posts or comments~\cite{avg260421446}. \textbf{Method and mechanism.} Each persona-conditioned agent repeatedly observes recent posts, interactions, its feed, and the social graph; a multimodal language model then selects posting, commenting, visual reply, liking, following, waiting, or other supported actions. For a visual reply, the selected target image is supplied as pixels to GPT-4o, which reasons over it and constructs a prompt for FLUX. The reply and resulting interactions become environment state available to subsequent decisions. Memory is episodic: recent activity is reloaded at each cycle, while no cross-session learned memory is accumulated. \textbf{Evaluation and results.} At the reported snapshot, 1,007 agents had produced 6,255 posts, 8,286 comments, and 4,684 image-bearing replies. The study identifies 498 visual chains of depth at least two, with mean depth 4.26 and maximum depth 60. Consecutive-image chain coherence is 0.702, compared with a random-pair null of 0.626; posts in chains average 13.0 recorded interactions versus 1.7 outside chains, although chain replies are included in the engagement count. \textbf{Evidence boundary.} Perceived images and social feedback change later actions, supporting L3 environment-conditioned behavior. AI-Gram is primarily a social-dynamics research platform rather than a method optimized against a user-specified visual artifact. Fixed personas, episodic memory, the weak coupling between reasoning and FLUX, and observational correlations prevent an L4 claim about persistent creative state or robust long-horizon artifact recovery.

\paragraph{InterleaveThinker.}
\textbf{Problem and setting.} InterleaveThinker addresses long-horizon interleaved generation, where a system must produce a sequence of mutually consistent images and text while preventing local visual feedback from overriding the global objective or propagating an error through later steps~\cite{avg260613679}. \textbf{Method and mechanism.} A Planner first converts the request into a global $N$-step plan. At step $i$, a frozen image generator conditions on the accepted image from step $i-1$. A Critic compares the pre- and post-action images with the current instruction, returns a binary judgment and an actionable prompt revision, and repeats generation until the step passes or a five-iteration limit is reached. The Planner and Critic are initialized by supervised fine-tuning, while a dual-reward GRPO stage trains the Critic to combine judgment accuracy with correction quality. \textbf{Evaluation and results.} With FLUX.2-klein-9B, the framework reaches 66.3 average on UEval, compared with 49.1 for Emu3.5; on CoMM, its style and entity-consistency scores are 9.3 and 9.2 for interleaved inputs. It raises the same FLUX backbone from 0.47 to 0.73 on WISE and from 13.3 to 28.9 on RISE-Bench. Increasing the per-step refinement budget from one to five raises the UEval average from 60.2 to 66.3. \textbf{Evidence boundary.} The accepted image sequence is retained across steps, visual diagnoses cause local recovery, and each step terminates from observed evidence, supporting L4 long-horizon control. The global plan is produced in advance, however, and the paper does not demonstrate rollback to an earlier accepted step or experience retained across requests. Performance remains bounded by the frozen renderer's generative prior, particularly for concepts absent from its training corpus.

\paragraph{AuDiffusion.}
\textbf{Problem and setting.} AuDiffusion addresses controllable text-to-image generation while also reducing the spatial and computational costs associated with globally attended diffusion backbones~\cite{avgdoi101007s40747025022111}. \textbf{Method and mechanism.} A prompt agent expands and weights semantic requirements, a layout agent matches those requirements to candidate ControlNet modules, and an executor runs an AuMamba latent diffusion model that combines selective state-space scans with windowed local attention. Their intermediate messages, selected controls, generated image, and quality score are stored on a shared blackboard. If semantic, layout, and realism scores fall below a threshold, the executor diagnoses count, attribute, or layout errors and revises prompt weights, the selected ControlNet, or generation parameters before another bounded attempt. \textbf{Evaluation and results.} On conditional ImageNet at $256\times256$, AuDiffusion-G reports FID 2.21, Inception Score 274.12, precision 0.85, and recall 0.59; at $512\times512$, it reports FID 3.02 and Inception Score 268.31. In zero-shot MS-COCO evaluation it obtains FID 12.31 and CLIP 0.3007. Removing the multi-agent controller increases the reported L2 reconstruction error from 0.0295 to 0.0325 and reduces DINO consistency from 0.8711 to 0.8261. \textbf{Evidence boundary.} The blackboard retains within-request state and low scores cause targeted re-execution, establishing L3 feedback adaptation. The evaluation principally measures the trained generator, and the controller ablation is narrower than a trajectory-level recovery study. Reported failures under conflicting semantic and structural priors include inaccurate interactions, boundary separation, and local geometric distortion; no persistent cross-request learning is demonstrated.

\paragraph{Multi-agent Collaborative Pathways for Chinese Traditional Architectural Image Generation.}
\textbf{Problem and setting.} This system generates images of Chinese traditional architecture from non-specialist requests while attempting to preserve historically specific form, color systems, decorative symbols, and cultural meaning~\cite{avgdoi101038s41598025181307}. \textbf{Method and mechanism.} A domain knowledge base encodes entities, architectural components, spatial rules, symbolism, and historical versions. Intent understanding and prompt-generation agents ground vague requests in this knowledge; a Stable Diffusion LoRA plus Canny ControlNet renders the image. An aesthetic and cultural-relevance agent returns scores and error flags to a workflow scheduler. Depending on the diagnosed cause, the scheduler revisits prompt generation, changes rendering parameters, or returns to intent interpretation, with a maximum iteration count preventing unbounded execution. \textbf{Evaluation and results.} Sixty participants, comprising 20 domain experts, 20 relevant students, and 20 general users, score anonymized outputs. The MAS prototype reports $4.65\pm0.25$ for architectural form, $4.72\pm0.20$ for color fidelity, and $4.71\pm0.22$ for decorative-symbol consistency, compared with $4.50\pm0.28$, $4.58\pm0.26$, and $4.61\pm0.29$ for the same generator driven by expert-authored prompts. Its average CLIP score is $0.89\pm0.03$, compared with $0.86\pm0.04$ for that baseline. \textbf{Evidence boundary.} Visual and cultural error reports determine targeted rerouting, providing L3 artifact-conditioned correction. The evidence comes from one cultural case study and an author-built knowledge base; CLIP against the optimized prompt is not a direct measure of historical truth, and the evaluation does not isolate the contribution of repeated correction from knowledge grounding and model specialization. A single pass averages about 166.7 seconds, with additional iteration cost.

\paragraph{MUSES.}
\textbf{Problem and setting.} MUSES addresses text-to-image prompts that specify three-dimensional object position, depth, orientation, count, and camera view, attributes that are weakly represented by conventional two-dimensional layout control~\cite{ding2025muses}. \textbf{Method and mechanism.} A layout manager retrieves related two-dimensional demonstrations and uses chain-of-thought reasoning to lift boxes into depth, orientation, and camera parameters. A model engineer follows a retrieve--search--generate decision tree to obtain 3D assets, aligns them to the camera using a fine-tuned CLIP classifier, and adds newly acquired assets to a model shop. An image artist assembles the scene in Blender and derives depth and Canny controls for final diffusion rendering. \textbf{Evaluation and results.} On T2I-CompBench, MUSES reports 0.8726 for color binding, 0.4756 for two-dimensional spatial relations, 0.4639 for three-dimensional spatial relations, and 0.7720 for numeracy. The authors also introduce 50-prompt T2I-3DisBench. Its average InternVL score is 0.6156, with a corresponding user score of 0.62 from 30 participants on 20 sampled prompts; SD v3 reaches 0.4206 and 0.36, respectively. \textbf{Evidence boundary.} The demonstrated generation path is L2 structured tool orchestration: planning and asset choice precede one final rendering path, and the benchmark-based choice among control scales is part of evaluation rather than an artifact-conditioned runtime correction policy. Adding assets to the shop suggests extensibility, but the paper does not test experience-driven cross-task policy improvement. Results on the author-created benchmark also depend on InternVL scoring and a small user-evaluation subset.

Multi-agent organization can allocate expertise and generate diverse candidates, but it is not by itself evidence of artifact-conditioned control. The stronger systems expose shared state and let evaluators redirect subsequent role actions; fixed role handoffs and internal model communication remain narrower forms of coordination.

\subsubsection{Collaborative Editing, Completion, and Restoration}
These systems divide image modification into complementary planning, execution, evaluation, and specialized restoration or completion roles. Their evidence concerns whether this division improves an evolving edited artifact.

\paragraph{CCA.}
\textbf{Problem and setting.} Collaborative Competitive Agents (CCA) addresses compound image edits that require several localized, stylistic, geometric, or preprocessing operations and that exceed the reliable scope of a single editor call~\cite{hang2024cca}. \textbf{Method and mechanism.} Two generator agents independently decompose the request and execute tools. A discriminator asks requirement-specific visual questions, combines their answers with aesthetic assessment, compares both candidates, and returns subtask-level feedback. Later rounds can change tool order, select another editor, or tune parameters; a best-candidate bank supports comparison across rounds, and the quality competitor stops when the current result satisfies the request or after five rounds. \textbf{Evaluation and results.} In the reported human comparison, CCA receives 47\% of votes for instruction alignment and 44\% for visual quality, compared with at most 21\% and 23\% for InstructPix2Pix, MagicBrush, or InstructDiffusion. Hierarchical tool descriptions reduce invalid plan formatting from above 20\% to below 10\%. On 20 requests, quality-based early stopping reduces average execution from 20 tool calls over five rounds to 12 calls over approximately three rounds. \textbf{Evidence boundary.} Visual feedback changes later plans and parameters, establishing L3 artifact-conditioned repair. The best-candidate bank is short-lived request state rather than durable long-horizon memory: the authors report that sequential subtasks cannot access earlier images and identify a memory module as future work. Experiments are dominated by qualitative and preference evidence, and the full GPT-4 configuration averages 11.5 minutes per request.

\paragraph{MAIR.}
\textbf{Problem and setting.} MAIR considers complex restoration where scene degradations, imaging degradations, and compression artifacts coexist and a user may request selective rather than maximal restoration~\cite{avg250309403}. \textbf{Method and mechanism.} Its scheduler first uses degradation perception, a three-stage restoration prior, textual experience summaries, and the user instruction to construct an overall plan. Specialist expert agents then make fine-grained observations of the current image, choose among registered tools, execute an action, and reflect before passing the result to the next expert. The registry permits tools to be added with descriptions rather than rebuilding the full orchestration interface. \textbf{Evaluation and results.} The study evaluates 1,440 synthetic MiO100 images and two real-world 100-image test sets. On the paired real-world set, MAIR records 21.67 PSNR, 0.7271 SSIM, 0.3244 LPIPS, and 115.61 FID, compared with AgenticIR's 19.14, 0.6574, 0.3841, and 120.66. Its average time and tool calls on the real-world sets are 35.42 seconds and 1.82, compared with 63.04 seconds and 5.15 for AgenticIR. Replacing intelligent expert selection with random selection lowers PSNR by 0.78 dB and MUSIQ by 10.05. \textbf{Evidence boundary.} Expert-level observations and reflection influence tool selection on the evolving image, which supports L3 correction. The three-stage prior constrains the plan in advance, the stored experience is offline guidance rather than deployed learning, and the framework does not demonstrate stateful branch recovery or cross-request adaptation.

\paragraph{CREA.}
\textbf{Problem and setting.} CREA studies creative image editing and generation, where the requested artifact must satisfy an object or concept while also exhibiting novelty, visual diversity, and user-specific creative direction. It evaluates both edits that preserve source structure and unconstrained generated images~\cite{venkatesh2025crea}. \textbf{Method and mechanism.} A Creative Director interprets the concept, a Prompt Architect produces and fuses contrastive prompts, a Generative Executor calls Flux or an editing-conditioned backbone, and a Critic scores six creativity principles. A Refinement Strategist maps weak criteria to prompt changes and starts another generation or edit. The current image and the creativity index are retained between rounds; optional user guidance can modify the evolving prompt, and execution stops when a threshold is met or after three refinement rounds. \textbf{Evaluation and results.} The evaluation uses 24 objects and 25 prompts, yielding 600 images per task for editing and generation. For editing, CREA obtains CLIP $0.417\pm0.030$, LPIPS $0.414\pm0.157$, Vendi $3.70\pm1.97$, DINO $0.744\pm0.185$, FID $294.19$, and KID $14.02$, improving over the listed baselines on the combined metrics. For generation, it obtains CLIP $0.360\pm0.052$, LPIPS $0.709\pm0.057$, Vendi $10.44\pm2.15$, FID $248.67$, and KID $5.94$. A 50-participant study reports creativity scores of $3.74$ for editing and $4.16$ for generation; the ablation increases LPIPS diversity from $0.302$ to $0.414$ as principles, contrastive prompting, and self-enhancement are added. \textbf{Evidence boundary.} Critic scores alter the next prompt and image, supporting L3 feedback adaptation with optional human intervention and bounded stopping. Creativity criteria are partly LLM-mediated and the base generator can hallucinate people or darken backgrounds; the paper does not report durable experience transfer, rollback to earlier image states, or long-horizon project control.

\paragraph{Talk2Image.}
\textbf{Problem and setting.} Talk2Image targets multi-turn conversational image generation and editing, where later instructions should accumulate with earlier scene content and affect only the requested regions. The artifact is an image state that evolves across dialogue turns, with generation, editing, visual question answering, and chat as related operations~\cite{avg250806916}. \textbf{Method and mechanism.} A multi-turn intention parser converts dialogue history into a structured cumulative specification. Specialized agents handle generation, editing, segmentation, visual questioning, and chat through a DAG of executable operations. A persistent blackboard stores intention history, agent outputs, and feedback records. After execution, multi-view feedback scores semantic alignment and visual consistency; a low score causes retry or refinement, and the process stops when the score stabilizes above a threshold or the retry limit is reached. \textbf{Evaluation and results.} On T2I-CompBench, Talk2Image reaches UniDet $0.2676$ and human spatial score $0.5240$, compared with $0.2543$ and $0.4956$ for DALL-E 3. On MagicBrush multi-turn editing, it reports $L2=0.0298$, CLIP-I $0.9198$, CLIP-T $0.3157$, DINO $0.8730$, and human score $0.8264$; removing the intention parser lowers CLIP-T to $0.3022$, removing collaboration lowers DINO to $0.8356$, and removing feedback lowers CLIP-I to $0.9011$. \textbf{Evidence boundary.} Persistent dialogue state, cross-turn artifact dependence, feedback-driven retries, and threshold-based termination support L4 stateful long-horizon control. The system is zero-shot and uses a fixed tool set with Qwen2.5-VL-7B as executor; an incomplete bear-removal case shows segmentation failure, and the paper does not demonstrate experience transfer across conversations or L5 self-improvement.

\paragraph{Multi-Agent Amodal Completion.}
\textbf{Problem and setting.} Multi-Agent Amodal Completion reconstructs the hidden portions of an occluded or image-boundary-truncated object and returns a layered RGBA asset. It is included as adjacent evidence for localized image synthesis and compositing~\cite{avg250917757}. \textbf{Method and mechanism.} An Occlusion Agent identifies front--back relations, a Segmentation Agent extracts visible-object and occluder masks, a Boundary Agent estimates canvas expansion, and a Description Agent hypothesizes fine-grained appearance for hidden parts. These decisions define one inpainting mask and semantic prompt before a FLUX-ControlNet Inpainting Agent performs a single synthesis pass. Cross- and self-attention maps are fused with the known visible mask to construct the alpha channel. \textbf{Evaluation and results.} On the VG component of the Ao et al. benchmark, the method reports CLIP $28.011$, LPIPS $0.218$, feature similarity $0.855$, and SSIM $0.839$; on COCO-A it reports $27.572$, $0.250$, $0.831$, and $0.820$. On custom Internet/Daily-Life data, LPIPS is $0.151/0.094$, feature similarity is $0.853/0.916$, and SSIM is $0.910/0.943$. A 30-participant study gives it $3.53$ for completeness, $3.87$ for coherence, $3.60$ for plausibility, and $3.67$ overall. Runtime is $25.51$ seconds on an A800, compared with $43.30$ for PD-MC and $51.82$ for OWAAC. \textbf{Evidence boundary.} The paper explicitly performs collaborative reasoning before one non-iterative synthesis pass. Generated pixels are not re-observed to alter the action sequence, so the demonstrated behavior is L2 orchestration rather than feedback adaptation. Metrics assess visible-region consistency because hidden ground truth is unavailable, and segmentation errors remain a primary failure source.

\paragraph{ImageEdit-R1.}
\textbf{Problem and setting.} ImageEdit-R1 addresses compound image-editing instructions that combine several actions, target objects, and intended outcomes, for which a direct editor may lose global intent or execute incompatible changes~\cite{avg260308059}. \textbf{Method and mechanism.} A Qwen2.5-VL-7B decomposition agent maps the request and source image into actions, subjects, and goals. A sequencing agent orders the resulting subrequests, and a diffusion-based editor executes the sequence. Reinforcement learning trains the decomposition agent with action, subject, and goal rewards; the editor itself is not modified. \textbf{Evaluation and results.} The paper evaluates 328 PSR requests, 1,000 sampled RealEdit cases, and 1,000 sampled UltraEdit cases using GPT-4o and Gemini-2.5 judges. Averaged over the three benchmarks, the framework raises FLUX.1-Kontext-dev from 7.21 to 8.23, Qwen-Image-Edit from 8.39 to 8.85, and Nano Banana from 8.32 to 8.66. A goal-aware reward gives an average 8.19, compared with 7.92 without goal-specific supervision; improvements level off after roughly 4,000 reinforcement-learning examples. \textbf{Evidence boundary.} ImageEdit-R1 is L2 learned decomposition and routing at deployment: the action sequence is prepared before the editor runs, and the completed or intermediate image is not used to select a new action. Reinforcement learning provides offline policy improvement rather than retained experience from completed user edits. Scores rely on LLM-as-a-judge evaluation, and the method inherits failures of the selected editor and of the initial decomposition.

\paragraph{CAMEO.}
\textbf{Problem and setting.} CAMEO studies conditional image editing under multiple semantic, structural, and contextual constraints. The evaluated artifacts include road scenes with inserted anomalies and human images whose pose must be switched while preserving anatomical plausibility and scene context~\cite{pu2026cameo}. \textbf{Method and mechanism.} A hierarchical orchestrator interprets the request, retrieves or constructs reference evidence, and generates a candidate edit. A Quality Critic evaluates the intermediate image under task-adaptive constraints and returns structured deviations. A Refinement Editor translates those deviations into a targeted correction, after which the candidate is re-evaluated. The loop stops when active constraints pass their thresholds or an iteration limit is reached; ablations remove reference grounding, quality control, or iterative refinement. \textbf{Evaluation and results.} The road-anomaly study uses 10,000 images sampled from BDD100K across 30 anomaly categories and 10 weather conditions; CAMEO wins roughly 20 percentage points more often than direct editing on average, with road-scene judge scores around $8.25$--$8.45$ versus $7.10$--$7.85$ across backbones and judges. The pose-switching benchmark contains 1,000 Pexels images and 10,000 edited samples; CAMEO generally improves over direct editing, with scores around $8.08$--$8.47$ versus $6.32$--$8.46$. A 10-annotator study, 100 pairs per annotator, gives net human preferences from $+4.3$ to $+34.3$ across the four baselines. \textbf{Evidence boundary.} Quality feedback changes the next edit and termination decision, supporting L3 artifact-conditioned correction. The evaluation relies heavily on four vision-language judges, the human study is small, and iterative coordination adds computation. The paper does not demonstrate persistent cross-request state, rollback, or transfer of learned editing skills, so it does not support L4 or L5.

\paragraph{APE.}
\textbf{Problem and setting.} APE addresses prompt underspecification in text-to-image generation and instruction-guided image editing. Its immediate artifact is an enhanced textual instruction passed to a frozen downstream image model; it is intended to preserve constraints, supply grounded details, and reduce omissions in compositional requests~\cite{avg260600204}. \textbf{Method and mechanism.} The single-agent SAPE variant and multi-agent MAPE variant use a router, field-specific rewriters, and a composer to decide which semantic fields require expansion. The editing version additionally receives the source image and can decide whether rewriting is necessary. Supervised fine-tuning and GRPO/GDPO use downstream image rewards such as PickScore, CLIPScore, HPSv2.1, and GPT-based editing judgments to train the enhancer, while inference produces one enhanced prompt before the renderer is called. \textbf{Evaluation and results.} On UniGenBench, MAPE with Qwen3-1.7B reaches short/long scores of $0.8334/0.8624$ for Qwen-Image-2512, $0.7716/0.8405$ for Z-Image-turbo, $0.7710/0.8275$ for FLUX.2-klein-4B, and $0.8268/0.8578$ for FLUX.2-klein-9B; the corresponding Qwen3 baselines are lower. On ImgEdit, MAPE with Qwen3-VL-4B obtains overall scores of $4.15$, $4.32$, and $4.01$ with FLUX.2-klein-4B, FLUX.2-klein-9B, and Qwen-Image-Edit, respectively, with gains concentrated in extract, style, and adjustment cases. \textbf{Evidence boundary.} The router's tool and role selection is completed before downstream rendering, and final image scores are used for training rather than to alter a deployed trajectory. APE therefore supports L2 tool and role assistance, not artifact-conditioned correction. Its reward is outcome-only, decomposition can introduce unnecessary edits for simple instructions, and results depend on the chosen image backbone and automated judges.

Collaboration is most useful when criticism can be assigned to a responsible edit or tool choice. Shared candidate banks, evaluator feedback, and repair roles support this connection, although fixed tool inventories and short iteration limits constrain recovery.

\subsection{Native and Unified Agentic Models for Image Creation}
Native and unified models place reasoning, understanding, generation, and sometimes tool use within a common multimodal policy. This organization is discussed separately from the feedback and workflow mechanisms analyzed in earlier groups.

\subsubsection{Native and Unified Models for Reasoning and Generation}
Native and unified models integrate image understanding, reasoning, and generation in one learned policy. The grouping concerns the implementation locus of these functions, not an automatic claim of high autonomy.

\paragraph{ImAgent.}
\textbf{Problem and setting.} ImAgent studies adaptive test-time scaling for image generation and editing, where vague prompts and stochastic rendering make a fixed refinement strategy inefficient~\cite{avg251111483}. \textbf{Method and mechanism.} Built on a unified multimodal model, a policy controller observes the current prompt, image, and action history and chooses among direct generation or editing, chain-of-thought prompt enhancement, image-conditioned prompt revision, local detail refinement, best-of-$N$ sampling, and STOP. Actions update the shared observation history; invalid action strings fall back to direct generation or editing. \textbf{Evaluation and results.} The paper evaluates three generation benchmarks and four editing benchmarks with Bagel and Janus-Pro-7B. On R2I-Bench, Bagel improves from 0.54 to 0.62 and Janus-Pro-7B from 0.47 to 0.51. On RISEBench, Bagel improves from 6.1 to 13.1, and on ImgEdit-Bench from 2.89 to 3.15. The generation trajectories average 4.73 steps for Bagel and 3.98 for Janus-Pro-7B under a five-step cap. \textbf{Evidence boundary.} The current image and history determine later actions and stopping, providing clear L3 feedback adaptation and several short-horizon L4 ingredients. We retain an L3 assignment because the reported horizon is capped at five actions and the system provides no evidence for authority management, extended task recovery, or cross-task state retention.

\paragraph{UniReason 1.0.}
\textbf{Problem and setting.} UniReason unifies text-to-image generation and image editing for requests that require implicit cultural, scientific, spatial, temporal, or logical knowledge. It aims to produce an initial grounded image and then repair visual discrepancies through editing-like refinement~\cite{avg260202437}. \textbf{Method and mechanism.} World-knowledge-enhanced textual reasoning expands an underspecified request into explicit semantic and spatial guidance. After initial synthesis, the model observes the image and reasoning trace, produces a structured self-reflection, and applies an editing-style refinement. Training data are constructed with a generator, Gemini-2.5-Pro verifier, Qwen-Image-Edit refinement teacher, and final comparative judge; the unified understanding and generation model is then trained on single-turn reasoning and interleaved refinement samples. \textbf{Evaluation and results.} WISE contains 1,000 knowledge-informed prompts, UniREditBench contains 2,700 editing cases, and KrisBench contains 1,267 knowledge-intensive edits. UniReason obtains WISE overall $0.78$, compared with $0.70$ for BAGEL reasoning and $0.75$ for UniCoT. It reaches KrisBench $68.23$ and UniREditBench $70.06$, versus $60.18/50.96$ for BAGEL and $68.00$ on KrisBench for UniCoT. On general benchmarks it reports GenEval $0.90$, DPGBench $86.21$, ImgEdit $4.06$, and GEdit-EN $6.94$. Adding two-stage training, reasoning, and refinement raises WISE from $0.52$ to $0.58$, $0.73$, and $0.78$. \textbf{Evidence boundary.} A generated image is re-observed and the reflection conditions a later editing action, supporting L3 feedback adaptation. The data-construction verifier and judge also filter supervision, but this is not runtime self-improvement. Refinement follows a trained pattern without persistent project state, rollback, authority management, or experience transfer across deployed tasks, so the evidence remains below L4 and L5.

\paragraph{VisionCreator.}
\textbf{Problem and setting.} VisionCreator addresses general image and video creation requests that require a model to interpret design intent, decompose a complex objective, and execute a long sequence of heterogeneous generation tools rather than follow a task-specific template~\cite{avg260302681}. \textbf{Method and mechanism.} The model unifies Understanding, Thinking, Planning, and Creation (UTPC) within one trainable policy. Its Progressive Specialization Training uses general reasoning data and 4,000 expert-filtered VisGenData trajectories, while Virtual Reinforcement Learning trains tool planning in VisGenEnv with plan, format, tool-execution, result, consistency, and trajectory-coherence rewards. At inference, textual instructions, intermediate tool feedback, and visual states enter a persistent multimodal trajectory; the data report a mean of 15 steps, with 64\% of trajectories exceeding 20 steps. \textbf{Evaluation and results.} VisGenBench contains 1,200 tasks, comprising 400 image and 800 video requests across ten evaluation dimensions. VisionCreator-8B obtains a VLM-evaluated overall score of 0.581 and a success rate of 0.925, compared with 0.577 and 0.863 for GPT-5; human evaluation gives VisionCreator-32B an overall score of 3.42, with image and video success rates of 99\% and 96\%, compared with 3.19 for GPT-5 and 3.01 for Gemini2.5-Pro under the reported protocol. \textbf{Evidence boundary.} The persistent trajectory and intermediate observations support L3 feedback adaptation and provide unusually strong long-horizon evidence. We do not assign L4 because the paper does not demonstrate authority management or systematic artifact-conditioned recovery and adaptive stopping across the evaluated tasks. Moreover, VRL explicitly replaces real visual feedback with structural proxy rewards, leaving a sim-to-real gap governed by tool fidelity and the correspondence between plan correctness and visual quality.



\paragraph{VisionCreator-R1.}
\textbf{Problem and setting.} VisionCreator-R1 studies error accumulation in single-image generation, multi-image workflows, and image editing, where plan-driven visual agents can execute a plausible tool sequence without diagnosing intermediate visual deviations~\cite{avg260308812}. \textbf{Method and mechanism.} It extends VisionCreator's UTPC trajectory to Understanding, Thinking, Planning, Creation, and Reflection (UTPCR). An Act--Reflect--Think--Act loop inspects intermediate images, identifies deviations from instructions and prior context, and can issue targeted edits or regeneration. Reflection--Plan Co-Optimization first learns reflection on lower-noise single-image tasks, then mixes reflection-strong and planning-strong trajectories in VCR-SFT before multi-task reinforcement learning on VCR-RL; reflection rewards are computed from query-specific visual checkpoints and are combined with plan, format, tool, and result rewards. \textbf{Evaluation and results.} On VCR-Bench, automatic scores for VisionCreator-R1 are 0.532 on single-image, 0.700 on multi-image, and 0.836 on image-to-image tasks, compared with 0.515, 0.649, and 0.816 for Gemini2.5-Pro. The multi-image traces attain a plan score of 0.9746 and 31.0\% good reflections. On GEdit-Bench, the reported overall score is 7.23, compared with 7.03 for the Qwen-Image-Fast executor, with semantic consistency increasing from 7.26 to 7.60. \textbf{Evidence boundary.} Intermediate visual diagnosis can directly trigger corrective creation actions, supporting L3 feedback adaptation across both short and multi-image trajectories. The paper also shows that reflection learned on single-image tasks does not transfer directly to multi-image reinforcement learning: downstream generation stochasticity and tool errors make reflection credit assignment high variance. The available evidence does not establish authority management, general rollback, adaptive stopping, or cross-task retained improvement, so the system is not assigned L4 or L5.

A unified architecture can reduce handoff loss between visual interpretation and generation, but behavioral evidence still depends on whether observations change later actions. Models trained for a single integrated pass remain distinct from controllers that repeatedly inspect and repair artifacts.

\subsubsection{Models for Grounded, Spatial, and Open-World Creation}
This group extends unified policies to world knowledge, spatial reasoning, and tool-mediated open-world generation. It concentrates on models whose primary contribution is grounded creation within an integrated multimodal policy.

\paragraph{Unify-Agent.}
\textbf{Problem and setting.} Unify-Agent targets images of public figures, intellectual properties, cultural concepts, landmarks, and other long-tail entities for which factual identity and contextual attributes cannot be recovered from the user prompt alone~\cite{avg260329620}. \textbf{Method and mechanism.} A unified multimodal model diagnoses missing knowledge, performs textual and visual search, and transforms retrieved evidence into a compact grounded recaption. The final synthesis stage conditions on that recaption and selected visual evidence. Training uses 143,000 annotated trajectories to supervise understanding, evidence acquisition, recaptioning, and rendering within the same model. \textbf{Evaluation and results.} On FactIP, the full system raises the vanilla BAGEL backbone's overall score from 50.9 to 73.2 and relevance from 44.9 to 72.4. It obtains 0.77 overall on WISE, compared with 0.70 for BAGEL with chain-of-thought prompting, and 4.08 on KiTTEN, compared with 3.04 for BAGEL. Removing image search reduces FactIP overall to 56.2, and removing recaptioning reduces it to 62.9. \textbf{Evidence boundary.} The paper describes the deployed method as a shallow one-pass workflow: retrieval and recaptioning precede one render, with no observation, reflection, or replanning over the generated artifact. It therefore supports L2 world-grounded tool use rather than a closed-loop level. Limited image context in the base model constrains more complex behaviors, and the strongest factual scores remain below leading closed systems.

\paragraph{Boogu-Image-0.1.}
\textbf{Problem and setting.} Boogu-Image-0.1 is an open image-generation and editing model family designed for complex, implicit, multilingual, and text-rich requirements under a constrained training and inference budget~\cite{avg260713125}. \textbf{Method and mechanism.} The core generator uses Qwen3-VL-8B as its instruction encoder and is trained through a data-curation and captioning pipeline. At inference time, a DeepSeek-V4-Flash agent interprets the request, performs reasoning-oriented prompt rewriting, and routes between Base and Turbo variants according to task complexity; the report also discusses best-of-$N$ sampling and reflection as test-time scaling options. The released Base, Turbo, and corresponding Thinking variants therefore combine model-level generation with an upstream request-understanding layer. \textbf{Evaluation and results.} A blind Boogu Arena comparison collects more than 4,000 pairwise votes across nine models. On Qwen-Image-Bench, Base-Thinking obtains overall scores of 53.57 on Chinese prompts and 53.73 on English prompts, leading the evaluated open-source systems while remaining below several proprietary models. On LongText-Bench, Turbo-Thinking obtains 0.971 overall, but the authors explicitly note that the metric favors character correctness and can miss dense-layout artifacts. On ImgEdit-Bench, Base-Edit-Thinking scores 4.64 overall versus 4.51 for the non-thinking variant. \textbf{Evidence boundary.} The documented runtime behavior establishes L2 request interpretation, prompt transformation, and model routing. The report names reflection but does not specify or evaluate a deployed loop in which diagnosis of the generated artifact changes subsequent action, so L3 is not established. Strong benchmark performance should also be separated from agency: the model still exhibits world-knowledge gaps, anatomical failures under multi-person interaction and occlusion, limited language coverage, and VAE reconstruction errors.

\paragraph{ToolArtist.}
\textbf{Problem and setting.} ToolArtist targets open-world image requests that require semantic reasoning, external knowledge, and native image generation within a single policy~\cite{zhao2026toolartist}. \textbf{Method and mechanism.} The post-trained unified multimodal model dynamically coordinates reasoning, search-tool invocation, and image generation. Its supervised trajectories are collected from a teacher with search and generation tools, and Reason-Act-Draw GRPO combines intent and quality rewards. \textbf{Evaluation and results.} On WISE, ToolArtist reports 0.79 overall, compared with 0.77 for Unify-Agent and GenSearcher-Qwen-Image in the reported table; on WorldGenBench-Humanities it attains a 22.10 average knowledge-checklist score, compared with 15.58 and 13.66. \textbf{Evidence boundary.} The unified policy supplies L2 tool-assisted, knowledge-grounded generation because the available evidence does not show a generated artifact being diagnosed and repaired in a later action. Training-time post-training is not cross-task continual improvement at deployment.

Grounding and unified tool use improve the information available to a generator. The current evidence is strongest for controlled tool selection and knowledge-sensitive synthesis; post-generation diagnosis, recovery, and long-term improvement require separate demonstration.

\subsection{Data, Evaluation, and Safety for Agentic Image Generation}
Data construction, learning, evaluation, and safety research supports the development and assessment of image agents. These studies determine how controllers are trained, what claims of reusable experience are supported, and how quality or safety failures can be measured.

\subsubsection{Data Construction, Policy Learning, and Reusable Experience}
This group studies how visual-agent data, policies, workflows, and experience stores are constructed or improved. The central question is what, if anything, persists beyond a single trajectory and improves a later task.

\paragraph{Gen-n-Val.}
\textbf{Problem and setting.} Gen-n-Val addresses scarce and imbalanced detection and instance-segmentation data, especially rare LVIS categories, while filtering synthetic foregrounds with incomplete masks, multiple objects, incorrect classes, or unsuitable backgrounds~\cite{huang2025gennval}. \textbf{Method and mechanism.} A TextGrad-optimized LLM constructs detailed prompts, Layer Diffusion generates transparent single-instance foregrounds whose alpha channels provide masks, and a VLM validator checks class identity, single-view composition, object integrity, and background conditions before accepted instances are composited into training scenes. \textbf{Evaluation and results.} The system generates 727,393 LVIS instances and raises Mask R-CNN mask mAP from 21.7 to 25.6 and rare-category mAP from 9.6 to 17.2. On COCO, the reported YOLO11m mask mAP increases from 39.8 to 42.9 and rare-category performance from 45.4 to 49.0. \textbf{Evidence boundary.} Validation governs sample admission but does not diagnose and repair the rejected image. The contribution concerns downstream detection and segmentation data augmentation, so it provides L2 validation evidence for image-data generation rather than a general image-creation loop.

\paragraph{SIDiffAgent.}
\textbf{Problem and setting.} SIDiffAgent targets text-to-image requests whose intent is ambiguous or whose first output has local defects, including incorrect object counts and anatomy artifacts~\cite{avg260202051}. \textbf{Method and mechanism.} The training-free controller uses Qwen-family models for intent analysis, prompt refinement, generation, visual scoring, and localized editing. If the evaluator score falls below a threshold, it produces targeted corrective guidance and permits up to two editing attempts. Successful trajectories and failure descriptions are embedded in a FAISS-indexed memory; the five nearest records are retrieved to guide later prompts and corrections. \textbf{Evaluation and results.} On GenAI-Bench and DrawBench, the paper reports an average VQA score of 0.884 on GenAI-Bench. Relative to the Qwen-Image baseline, the basic workflow improves the reported score by 4.22\%, the complete system by 12.54\%, and the second episode using first-episode memory by 16.77\%. It also reports relative gains over Imagen3, Stable Diffusion 3.5, and T2I-Copilot. \textbf{Evidence boundary.} The score-triggered local edit establishes L3 output-conditioned correction, while the episode-to-episode memory experiment supplies L5 evidence for retained, reusable trajectory guidance. Results rely on the Qwen model family and model-based VQA scoring, memory retrieval can carry inappropriate guidance across tasks, and the bounded two-attempt policy leaves extended recovery untested.



\paragraph{Socratic-Geo.}
\textbf{Problem and setting.} Socratic-Geo addresses the shortage of executable, geometrically valid image--text--solution data for multimodal reasoning and diagram generation~\cite{avg260203414}. \textbf{Method and mechanism.} A Solver's failed reasoning attempts trigger a Teacher to diagnose the missing geometric constraint, modify parameterized Python drawing code, execute it, and apply two checks: RePI verifies visual validity and Reflect verifies solvability against the reference answer. Accepted triplets enlarge the next Solver curriculum. In parallel, the Teacher converts validated programs into drawing instructions used to fine-tune a diffusion Generator; the paper explicitly treats this Generator as an offline byproduct of the Teacher--Solver loop. \textbf{Evaluation and results.} Socratic-Generator obtains a strict score of 6.0 and a relaxed score of 42.4 on GenExam-Math, compared with 0.0 and 18.9 for its Qwen-Image base model and 8.0 strict for GPT-Image-1. For reasoning, the qualified-data variant reaches 40.33 on MathVerse using 0.4K samples, compared with 37.09 without qualification using 1.3K samples. Extensions raise ChartQA from 87.3 to 91.2 and UIFlow2Code from 75.9 to 81.5. \textbf{Evidence boundary.} At the framework level, verified failures alter later curriculum construction and repeatedly update reusable Solver and Generator models, with held-out and cross-domain gains; this supports L5 training-time self-improvement under the survey's criterion. The deployed image Generator remains a one-shot model and does not inspect or repair its own output. Thus the result is adjacent evidence for self-evolving visual-data production, not evidence of an L5 online diagram-creation agent. Training also uses 32 A100 GPUs and GPT-5-based automatic judging for the generation benchmark.

\paragraph{PaAgent.}
\textbf{Problem and setting.} PaAgent addresses image restoration across isolated and combined degradations, including snow, noise, haze, rain, low light, and blur. Despite its name, the studied output is a general restored image; the controller is intended to select a succession of restoration experts that improves both subjective appearance and objective image quality~\cite{avg260317055}. \textbf{Method and mechanism.} A Qwen-based perception module reads the current image, no-reference quality scores, and completed tasks. Retrieval-augmented reasoning queries a tool portrait bank whose records associate degradation descriptions, an applied tool, and a subjective evaluation. The selected expert updates the image; its outcome is summarized and added to the bank. Subjective-objective reinforcement learning trains the next-task policy using multimodal evaluations and seven no-reference IQA metrics, and the runtime repeatedly evaluates the intermediate result before choosing another task or terminating. \textbf{Evaluation and results.} Training uses 15,013 images and the test suite contains 16,230 pairs spanning eight benchmarks. PaAgent reports 37.95 PSNR and 0.98 SSIM on CSD desnowing, and at noise level 50 on BSD68 it reports 27.03 and 0.79. On SOTS indoor dehazing, it obtains 40.65 PSNR and 0.99 SSIM; on Rain13K Rain100H it obtains 29.57 and 0.90. In the LOL Blur task, it reaches 20.82 PSNR and 0.83 SSIM. On four CDD-11 mixed-degradation settings, the full method reaches 27.83/0.91, 27.02/0.86, 26.46/0.87, and 25.95/0.86 PSNR/SSIM, whereas the supervised-only ablation reaches 25.31/0.86, 24.84/0.78, 24.10/0.75, and 23.59/0.71. \textbf{Evidence boundary.} Observation of intermediate images, executed-task history, adaptive continuation, and termination establish L4 stateful restoration control. The tool portrait bank persists newly summarized interactions, but the current evaluation does not isolate whether experience accumulated from completed deployment tasks improves performance on a separate later task set. The GRPO policy update is performed in training, so the available evidence does not establish L5 continual improvement; results also depend on proprietary Qwen evaluation and retrieval services.

\paragraph{ScaleEdit-12M.}
\textbf{Problem and setting.} ScaleEditor addresses construction of large, diverse, open-source image-editing data without relying on costly commercial generation APIs. Its primary artifact is ScaleEdit-12M, a corpus of instruction--source--edited-image triplets rather than an interactive editor for end users~\cite{avg260320644}. \textbf{Method and mechanism.} Source-image expansion combines open datasets and web retrieval with validity and duplicate filtering. A task router assigns images to specialized editing agents spanning 23 task families, including global, object, attribute, text-aware, knowledge-infused, and compositional edits. Task-aware verification scores instruction following, edit consistency, and generation quality; only pairs with a perfect instruction score and adequate consistency and quality are retained. The accepted data then supervise fixed downstream image editors. \textbf{Evaluation and results.} The resulting corpus contains 12 million pairs, of which $85.3\%$ receive the highest minimum score across all three verification dimensions before final filtering. UniWorld-V1 trained on ScaleEdit reaches $6.55$ on GEdit-EN-Full versus $4.85$ for its baseline, while Bagel reaches $7.17$ and $3.45$ on GEdit-EN-Full and ImgEdit. On RISEBench/KRIS-Bench, ScaleEdit-trained UniWorld-V1 reaches $5.55/56.60$ and Bagel reaches $7.50/63.58$. A 20-expert, 1,000-sample verification study reports Qwen2.5-VL-72B agreement with humans of $0.78$, $0.67$, and $0.78$ for instruction, consistency, and quality. \textbf{Evidence boundary.} Routing and verification provide L2 orchestration for dataset production: failed samples are filtered, but the system does not demonstrate that their artifact evidence causes a targeted regeneration trajectory. Downstream fine-tuning gains measure data utility rather than controller autonomy. The work is therefore adjacent evidence for scalable visual-data construction; generator quality caps the corpus, expert training over 23 tasks is expensive, and multi-turn editing is left for future work.



\paragraph{GenEvolve.}
\textbf{Problem and setting.} GenEvolve studies open-ended image generation in which an agent must decide when to retrieve facts and reference images, activate generation knowledge, and translate those resources into an executable prompt--reference program~\cite{chen2026genevolve}. \textbf{Method and mechanism.} Each attempt is a multimodal tool trajectory ending in a prompt and an ordered reference set. After supervised trajectory tuning, group-relative policy optimization samples several attempts per request and scores both the generated image and the program. Tool-Orchestrated Visual Experience Distillation compares the best and worst trajectories, abstracts differences into search, knowledge, reference, prompt, and failure-avoidance slots, conditions a privileged teacher on retrieved experience, and distills the resulting token-level preferences into the student policy. \textbf{Evaluation and results.} On GenEvolve-Bench, the open-generator configuration reaches a KScore of 0.3663, compared with 0.3493 for Gen-Searcher using the same Qwen-Image-Edit executor; pairing GenEvolve with Nano Banana Pro reaches 0.5739, compared with 0.5298 for the raw generator. On the external WISE benchmark, GenEvolve obtains an overall WiScore of 0.82, compared with 0.78 for Mind-Brush and 0.77 for Gen-Searcher. An ablation progresses from 0.3480 for SFT, to 0.3548 for SFT plus GRPO without visual experience, and to 0.3663 for the full method. \textbf{Evidence boundary.} Experience extracted from earlier training trajectories changes the reusable student policy and improves held-out requests, satisfying the survey's L5 cross-task improvement criterion. This is post-training self-evolution rather than online continual adaptation: the deployed student has no runtime visual-experience memory. Results also rely on proprietary multimodal judges and, in the strongest setting, a proprietary generator, so the policy gains should be separated from executor quality.

\paragraph{EvoIR-Agent.}
\textbf{Problem and setting.} EvoIR-Agent restores images with coupled degradations while reducing the trial-and-error cost of selecting both a restoration order and a specialized tool. It treats the restored image, the selected tool path, and the target preference for fidelity or perceptual quality as information that can be reused in subsequent cases~\cite{avg260522208}. \textbf{Method and mechanism.} The runtime follows perception, planning, execution, reflection, and rollback. Its hierarchical experience pool stores three forms of guidance: high-level textual insights, degradation-type mappings, and retrieved fine-grained degradation-pattern records. After a batch of restoration records accumulates, an experience-evolution procedure adds, replaces, updates, or deletes entries, then stabilizes their priorities from quality rankings. Thus later planning and rollback can retrieve guidance that reflects previous restoration outcomes. \textbf{Evaluation and results.} The evaluation includes 1,440 synthetic MiO100 images and two 50-pair FoundIR settings. Under the AgenticIR-compatible tool setting, EvoIR-Agent obtains Group A/B/C PSNR of 23.49/23.41/21.52, compared with 21.04/20.55/18.82 for AgenticIR. On FoundIR low-light plus noise, it reaches 17.81 PSNR and 0.7093 SSIM with 4.73 tool invocations, compared with 16.77, 0.6410, and 15.91 for 4KAgent. After two experience-evolution rounds on Group B, the reported unified quality index rises from 0.0761 to 0.9824 while total rollback decreases from 2.50 to 1.29. For the out-of-distribution FoundIR test, learning a fresh pool from 20 images and evaluating on the remaining 30 yields 24.50 PSNR and 0.8272 SSIM for low-light plus noise, compared with 17.93 and 0.7064 when reusing the original zero-shot pool. \textbf{Evidence boundary.} Intermediate reflection and rollback provide L4 recovery, while the experience pool is explicitly updated and transferred to later held-out restoration cases, supporting L5 cross-task mechanism improvement. This is batch-wise adaptation using benchmark data and quality signals, not autonomous unrestricted deployment; transferred prior experience can harm an out-of-distribution setting, as the reported zero-shot low-light-plus-noise result shows.

\paragraph{OctoT2I.}
\textbf{Problem and setting.} OctoT2I addresses selection among text-to-image models with different quality, latency, and energy profiles, where fixed routing priors cannot track each model's capability frontier~\cite{jiang2026octot2i}. \textbf{Method and mechanism.} During inference, a stateful multi-round router consults long-term tool knowledge and task-local execution memory, chooses a generator, evaluates prompt--image alignment, and stores the score for the next decision. Its self-evolving Propose--Solve--Evaluate--Learn loop defines conceptual dimensions, generates exploration prompts, measures each tool's Pass@1, prunes mastered or unproductive combinations, and updates both empirical records and semantic capability profiles. \textbf{Evaluation and results.} With five generation tools, OctoT2I scores 0.96 on GenEval and 0.6618 on T2I-CompBench++, compared with 0.93 and 0.6332 for Flow-GRPO. Average GenEval inference time is 10.02 seconds, versus 19.07 seconds for Flow-GRPO. Replacing learned knowledge with expert-written priors lowers GenEval from 0.96 to 0.93; on WISE, adding one newly explored tool raises the average from 0.54 to 0.61 using 60 exploration prompts, and adding a second raises it to 0.71. \textbf{Evidence boundary.} Current-task scores alter later routing, satisfying L3, while self-interaction updates long-term knowledge that improves later prompts and integrates unseen tools, providing L5 cross-task improvement evidence under this survey's criterion. This improvement updates a router knowledge base rather than generator weights; evaluation uses model-based alignment scorers, the default decision horizon is four rounds, and the system does not demonstrate L4-style authority management or extended project recovery.

\paragraph{A Task-Driven and Quality-Assured Agent Framework for SAR Data Generation.}
\textbf{Problem and setting.} SAGA generates synthetic-aperture-radar augmentation data for downstream recognition tasks, where a useful image-generation workflow must respect dataset schemas and assess whether its outputs meet SAR-specific quality constraints~\cite{avg260628896}. It is adjacent evidence because the artifact is specialized training data rather than an end-user visual creation. \textbf{Method and mechanism.} The system profiles a dataset against a schema, ranks alternative generation plans, executes a recipe DAG, and uses observer checks with bounded repair. A policy memory stores planning and verification evidence so later requests can reuse validated recipes. \textbf{Evaluation and results.} Experiments use 210,105 SAR images, 35 schema cases, and 80 planning requests. SAGA makes 34 of 35 schema decisions correctly, accepts 95\% of valid cases and rejects all invalid cases; plan selection reaches 92.5\% Top-1 and 97.5\% Top-3 with no invalid skill selections. The observer attains 90.0\% recall and 91.8\% F1, and the full system improves eight downstream task groups by 7.4 points on average over no augmentation. Removing memory lowers the reported gain from 7.25 to 6.33 and the composite score from 90.4 to 84.9. \textbf{Evidence boundary.} Cross-run policy memory provides adjacent L5-style evidence for reusable data-construction knowledge, while the observer and bounded repair provide artifact-conditioned correction. The evidence is SAR-specific, relies on private data and task-specific quality criteria, and does not show that the same controller transfers to general image creation.

Evaluation, methodological, and safety studies make the corresponding capability or failure mode observable without being counted as another generation architecture.

\paragraph{Self-Evolving Agentic Image Restoration via Deliberate Planning and Intuitive Execution.}
\textbf{Problem and setting.} SEAR addresses restoration of images with coupled degradations, for which the quality of a restoration depends on a sequence of specialized operations and early mistakes can remove information needed by later operations~\cite{avg260628971}. \textbf{Method and mechanism.} A deliberate planner uses language-model scheduling and pruning-aware Monte Carlo tree search to explore restoration paths. A hybrid no-reference reward ranks candidate paths, with a multimodal tournament used as an additional safeguard. High-reward trajectories are distilled into episodic memory indexed by degradation-aware state fingerprints; when a similar state recurs, an intuitive executor retrieves the stored strategy rather than searching again. \textbf{Evaluation and results.} The paper evaluates 1,440 MiO100-derived synthetic images covering 16 mixed degradations and a 100-pair real-world set assembled from six restoration benchmarks. On the synthetic Group B setting, SEAR reports 22.1332 PSNR, 0.7251 SSIM, and 0.2890 LPIPS, compared with 21.7160 PSNR for AgenticIR. On the real-world suite, it reports 24.4078 PSNR, 0.7425 SSIM, 0.3371 LPIPS, 0.4519 CLIP-IQA, and 54.6686 MUSIQ. Removing episodic memory increases the average number of tool calls from 8.15 to 16.75. \textbf{Evidence boundary.} Artifact-conditioned search and reuse support stateful long-horizon restoration, and the retrieved trajectories are evaluated as reusable knowledge for later cases, providing L5 evidence under the survey criterion. The claimed improvement remains conditional on the state-fingerprint retrieval scheme, the fixed restoration-tool library, and the calibration of no-reference and multimodal evaluators; the experiments do not establish safe generalization to arbitrary degradation distributions.

\paragraph{DataEvolver.}
\textbf{Problem and setting.} DataEvolver concerns construction of training data for text-rich image generation, where one-time crawl--filter pipelines discard rejected samples even though OCR, semantic, layout, and quality failures can guide later collection rounds~\cite{avg260631537}. \textbf{Method and mechanism.} A Retriever gathers topic-aware candidates; a Verifier accepts or rejects them using perceptual quality, OCR, semantic consistency, and duplicate checks; a Critic summarizes pass rates and rejection causes into semantic feedback; and a Generator fills under-covered regions with targeted synthesis. The resulting feedback memory updates later retrieval queries and generation prompts while deduplicating near-identical experience entries. \textbf{Evaluation and results.} Under a matched 0.75M-data budget, PixArt-$\alpha$ trained on DataEvolver data reaches OCR-F1 of 8.45 on TextScenesHQ and 9.08 on LongTextBench, versus 4.56 and 6.71 for the strongest reported baseline. For Show-o2, the corresponding values are 0.45 and 0.44, versus 0.19 and 0.27. At the 0.1M scale, removing the Critic reduces OCR-F1 from 1.78 to 1.01 on TextScenesHQ and from 2.16 to 0.90 on LongTextBench; removing targeted generation yields 1.40 and 1.37. \textbf{Evidence boundary.} Round-level verification is retained and changes later data-construction actions, establishing L5 evidence for a reusable construction policy. This is evidence about improving the data supply for a downstream renderer, not an end-to-end controller that repairs its own generated image during user-facing inference. Reported scaling is a trend rather than a fitted law, and noisy OCR, semantic scoring, or image-quality estimates can propagate into later policy updates.

\paragraph{COMFYCLAW.}
\textbf{Problem and setting.} COMFYCLAW targets recurring image-generation requests implemented as ComfyUI workflows, where an agent must construct a valid typed graph, meet node and model constraints, diagnose visual failures, and avoid rediscovering the same workflow repair in later tasks~\cite{avg260701709}. \textbf{Method and mechanism.} The harness exposes stage-gated tools for typed directed-acyclic-graph editing, parameter setting, execution, and rollback of invalid edits. A region-level VLM verifier converts visual defects into repair suggestions. Successful trajectories, execution errors, and verifier feedback are distilled into progressively disclosed skills; a proposed skill is validated on held-out tasks before it is committed to the library and becomes available to later workflows. \textbf{Evaluation and results.} Experiments span GenEval2, DPG-Bench, OneIG-EN, and OneIG-ZH, two image backbones, and three agent models. For Claude-Sonnet-4.5 with Z-Image-Turbo, COMFYCLAW raises the four-benchmark average from 67.94 to 77.78; with LongCat it raises the average from 67.08 to 75.52. The paper reports a four-point advantage over a harness-only verifier baseline and a ten-point advantage over no-refinement control, plus a 2,400-image human study. It retains 318 evolved skills, which account for roughly half of later skill invocations. \textbf{Evidence boundary.} Validated skills are created from prior trajectories and reused in later tasks, supporting L5 cross-task mechanism improvement. The evidence concerns ComfyUI graph construction rather than unrestricted visual creation, and all four main scores are judged by Qwen3-VL-8B-Instruct; skill quality remains sensitive to verifier error, benchmark coverage, graph-complexity constraints, and compatibility of the available nodes and models.

\paragraph{Search Beyond What Can Be Taught.}
\textbf{Problem and setting.} This work studies world-knowledge-intensive requests for which indiscriminate web or image search can corrupt concepts, copy irrelevant visual structure, or retrieve information that the current generator already represents adequately~\cite{wang2026searchgen}. \textbf{Method and mechanism.} A noise-resistant reasoner applies a Gate--Filter--Integrate protocol: it decides whether and in which modality to search, selects evidence that fills a specific knowledge gap, and converts the retained content into a grounded generation specification. Co-training then moves the generator--reasoner knowledge boundary in two stages. Online DPO teaches the generator stable knowledge and robustness to imperfect references; rejection-sampling fine-tuning recalibrates the reasoner to search only for knowledge that remains outside the updated generator. \textbf{Evaluation and results.} On the held-out SearchGen-Bench evaluation, the Klein-4B configuration improves monotonically from 26.4 after reasoner SFT, to 29.2 after generator DPO, and to 31.8 after reasoner recalibration, slightly above the 31.2 frontier-oracle result on the same generator. The corresponding Bagel progression is 23.4, 24.7, and 26.8. Pairing the recalibrated reasoner with the original rather than the DPO generator reduces Klein-4B performance from 31.8 to 26.8, supporting generator-specific policy calibration. \textbf{Evidence boundary.} Generator and reasoner parameters are updated from visual outcomes, reused on unseen prompts, and validated across two generator architectures, satisfying L5 cross-task self-improvement in this survey. The experiment contains one offline DPO pass followed by one offline reasoner-training pass; the authors identify multi-round co-training as future work. It therefore does not demonstrate recursive online self-evolution during deployment, and its reward signal depends on a Gemini-3-Flash judge despite reported human correlation.

The strongest continual-learning claims retain and reuse skills, toolpaths, or experience on held-out tasks. Offline reinforcement learning, self-generated data, or a self-improvement label alone do not establish that a deployed controller accumulates reusable experience.

\subsubsection{Evaluation and Safety}
Evaluation and safety studies provide the measurements and adversarial evidence needed to assess image-agent behavior. They are not treated as image-creation controllers merely because they use planning, search, or multiple roles.

\paragraph{A Unified Agentic Framework for Evaluating Conditional Image Generation.}
\textbf{Problem and setting.} CIGEval evaluates conditional image generation, where a score must account for source-image preservation, multiple concepts, and diverse controls rather than prompt alignment alone~\cite{avg250407046}. \textbf{Method and mechanism.} Its evaluator decomposes the assessment task, selects visual-analysis tools, and produces a fine-grained explanation alongside a score. The authors also synthesize evaluation trajectories to fine-tune smaller multimodal models. \textbf{Evaluation and results.} Across seven conditional-generation tasks, the GPT-4o version reaches a 0.4625 correlation with human assessments, close to the reported 0.47 inter-annotator correlation. A 7B open-source evaluator is fine-tuned using 2.3K filtered evaluation trajectories and is reported to exceed prior GPT-4o-based evaluation baselines. \textbf{Evidence boundary.} This is an evaluation architecture rather than a visual-creation controller, so no L-level is assigned. It supplies evidence about verifier design and explainability, while its reported agreement still depends on the human annotation protocol, the selected toolbox, and model-mediated trajectory synthesis.

\paragraph{EdiVal-Agent.}
\textbf{Problem and setting.} EdiVal-Agent evaluates multi-turn and compound image editing when each instruction can modify one object while later edits must preserve satisfied constraints and unaffected regions~\cite{avg250913399}. \textbf{Method and mechanism.} The evaluator converts edit requests and image history into object-centered criteria, assesses instruction following, object and background consistency, and visual quality at each turn, and aggregates the evidence into task-level results. Its benchmark annotations cover 724 objects in 13 categories. \textbf{Evaluation and results.} The paper reports per-instruction and per-turn success rates for multi-turn and complex-edit modes, together with DINOv3 and pixel-based consistency measures and human-preference scores. The breakdown makes clear that scores differ substantially across edit type and turn, exposing accumulated preservation failures that a final-image-only assessment conceals. \textbf{Evidence boundary.} EdiVal-Agent is an evaluation instrument rather than an autonomous editor, so it receives no L-level. Its measurements depend on automated object decomposition, evaluator judgments, image resizing, and the selected GEdit-Bench distribution; it provides a concrete source of diagnostic evidence for L3--L4 editing systems.

\paragraph{Blueprint-Bench.}
\textbf{Problem and setting.} Blueprint-Bench tests whether systems can transform interior photographs into an accurate two-dimensional apartment floor plan, making spatial connectivity and relative room size explicit visual constraints~\cite{avg250925229}. \textbf{Method and mechanism.} The benchmark contains 50 apartments with about 20 interior images each. It extracts a room-connectivity graph from each produced plan and computes a weighted similarity from edge overlap, degree correlation, graph density, room count, door count, and door orientation. \textbf{Evaluation and results.} It compares leading language models, image generators, and computer-use agents against random and human baselines. Most models are at or below the random baseline; agentic iterative refinement does not yield statistically meaningful improvement, and all tested systems remain substantially below the human result on the 12-apartment human-comparison subset. \textbf{Evidence boundary.} Blueprint-Bench is a diagnostic benchmark, not a generation method or evidence of an autonomy level. Its score omits room type and shape, and strict output formatting can confound spatial reasoning with instruction following; nevertheless, it identifies an empirical gap that a closed-loop visual controller must address.

\paragraph{Value-Aligned Prompt Moderation.}
\textbf{Problem and setting.} VALOR addresses unsafe, ambiguous, and value-sensitive requests to text-to-image systems, where a prompt-only filter may miss indirect intent and an unconditional refusal may remove legitimate creative content~\cite{avg251111693}. \textbf{Method and mechanism.} A zero-shot controller combines layered prompt analysis, value reasoning, and intention disambiguation. It selectively rewrites a problematic request, generates an image, and applies a post-generation safety check; an unsafe output can trigger a style-guided regeneration that retains the permitted semantic content. \textbf{Evaluation and results.} Experiments cover adversarial, ambiguous, and value-sensitive prompts. The paper reports reductions in unsafe outputs of up to 100.00\% while retaining prompt usefulness and creative diversity under its evaluation protocol. \textbf{Evidence boundary.} The image-level safety check can change the next generation attempt, giving bounded L3 correction for safety moderation. VALOR is not a general visual creator and does not retain task experience, so it does not establish L5 improvement; its safety claims remain sensitive to the threat set, value rubric, and post-hoc safety classifier.

\paragraph{OrchJail.}
\textbf{Problem and setting.} OrchJail studies safety failures in tool-calling text-to-image agents, focusing on attacks that distribute a harmful objective across individually acceptable planning and tool-execution steps~\cite{avg260507414}. \textbf{Method and mechanism.} Its orchestration-guided fuzzer mutates requests, plans, and tool-call sequences, observes safety decisions and generated outcomes, and searches for compositions that bypass guards at the controller--tool interface. \textbf{Evaluation and results.} The paper evaluates the resulting attack trajectories against tool-calling image-generation agents and reports successful bypasses that are missed by safeguards applied to individual prompts or tool calls. The study localizes failures to orchestration and cross-step semantic composition rather than only to the underlying image model. \textbf{Evidence boundary.} OrchJail is a red-team evaluation, not an image-creation controller and therefore receives no L-level. Its results demonstrate an attack surface and depend on the tested tools, policies, and threat model; they should not be interpreted as a quality or autonomy result for the affected generators.

\paragraph{Whispers in the Noise.}
\textbf{Problem and setting.} Whispers in the Noise examines whether concept-erased diffusion models can be induced to recover suppressed concepts, including safety-sensitive content, through interventions in the denoising trajectory~\cite{avg260518150}. \textbf{Method and mechanism.} The proposed ConceptAgent is a training-free black-box multi-agent procedure that constructs surrogate visual priors, reasons about the generation trajectory, and alters intermediate denoising states through surrogate-guided noise initialization. \textbf{Evaluation and results.} Across diffusion backbones and concept-erasure methods, the study evaluates concept-awakening accuracy and CLIP-based generation fidelity. It reports competitive or stronger scores than the compared training-based and zero-shot baselines, and demonstrates recovery of both ordinary and safety-sensitive erased concepts. \textbf{Evidence boundary.} This is safety evidence rather than a creative-generation architecture, and it receives no L-level. The method executes a single inference pass without online optimization or retained cross-task learning; its conclusions rely on surrogate quality, the selected erased concepts, and the benchmark classifiers used to recognize them.


\paragraph{RedEdit.}
\textbf{Problem and setting.} RedEdit tests whether image-safety classifiers remain reliable when an initially benign photograph is transformed through a sequence of ordinary editing operations~\cite{avg260606140}. \textbf{Method and mechanism.} The red-team agent formulates editing as a search tree, uses Monte Carlo tree search to select image edits, queries the target classifier as feedback, and expands paths that increase the target unsafe score while controlling visual plausibility. \textbf{Evaluation and results.} Experiments measure classifier vulnerabilities under iterative photo editing and show that the search procedure discovers adversarial edit trajectories that evade static safety judgments more effectively than non-search baselines. The analysis attributes failures to the classifier's limited robustness to composed semantic changes. \textbf{Evidence boundary.} This is a controlled safety attack rather than a constructive AVG method, so it has no autonomy-level assignment. Its evidence is useful for evaluating verifier robustness, but attack success is conditioned on the target classifier, edit operators, query budget, and safety definition.

\paragraph{Does AI Understand Imaging?.}
\textbf{Problem and setting.} ImagingBench tests whether multimodal agents can reason about computational imaging, including image formation, inverse reconstruction, sensing, optics, and calibration, rather than produce visually plausible but physically inconsistent outputs~\cite{avg260707189}. \textbf{Method and mechanism.} The benchmark defines 20 subtasks in five categories and three protocols: fixed expert-guided reconstruction, planner-guided reconstruction, and forward-system simulation for consistency checking. It evaluates proprietary and open image-centric multimodal systems alongside task-specific non-agentic baselines using raw task metrics and normalized aggregates. \textbf{Evaluation and results.} Agentic systems remain consistently weaker than specialized methods, with particularly large gaps in lensless imaging, event reconstruction, time-of-flight imaging, and holography. Planner guidance produces only modest and inconsistent gains over the fixed expert-prompt protocol, while visually plausible outputs often have poor reference-based fidelity. \textbf{Evidence boundary.} ImagingBench is a benchmark rather than an image-generation architecture, so no L-level is assigned. Its evidence is nevertheless directly relevant to AVG because it distinguishes semantic plausibility from physical and reference-grounded correctness; conclusions are limited to the implemented task suite and available model interfaces.

These studies broaden assessment beyond final appearance to object-level edit correctness, spatial reasoning, physical fidelity, safety moderation, and verifier robustness. Their diagnostic outputs become evidence for AVG only when a separate controller uses them to choose a corrective action.

Across image tasks, corrective precision depends on whether the system retains an addressable representation of the artifact and connects verification evidence to a later action. Masks, object tables, edit histories, preference records, degradation estimates, and executable workflows can support localized revision; undifferentiated scores often lead to global resampling. The evidence remains uneven for verifier calibration, causal diagnosis, rollback, budget-aware stopping, and persistent improvement beyond one task.

\FloatBarrier
